\def\ChapterCode{TRA}
% \begin{savequote}
% \includegraphics[width=3.5cm]{./images/wikipedia-pd/Battlefield_wounds_300px.jpg}
% 
% {\tiny\itshape Hans von Gersdorff, Feldbuch der Wundarznei, 1517.}
% \end{savequote}
\chapter
[\CO{[TRA]} Trauma]
{Trauma}
\label{chp:TRA}
% \AasRsN{RS~VII}

\ChapterIntro
{%
 Als \T{Trauma}  bezeichnet man eine
\begin{inparaenum}
  \item Schädigung,
  \item Verletzung, oder
  \item Wunde,
\end{inparaenum}
die durch eine
Gewalteinwirkung physikalischer (mechanischer,  thermischer)
oder chemischer Art verursacht wurde.
}
{%
\begin{description}
  \item[Maintainer:] Michael Motal
  \item[Autoren:] Diverse
  \item[Reviewer:] Standard-Reviewprozess
  \item[Version:] {\DocVersion} ({\DocVersionInfo})
  \item[SHA1:] {\DocGitCommit}
\end{description}

{\TxTeilkapitelAASS}
}


\clearpage


Die Untersuchung des traumatologischen Patienten
(Traumacheck) wird unter \prettyref{topic:traumacheck}
erklärt.

Auch die Psyche kann als Resultat einer erheblichen
seelischen Belastung bzw. unzureichender
Bewältigungsmöglichkeiten ein Trauma erleiden. Dies zeigt
sich dann z.B. als Belastungsstörung oder
Persönlichkeitsänderung. Vgl. Kapitel \CO{[PSY]}
(\prettyref{chp:PSY}).





%%%%%%%%%%%%%%%%%%%%%%%%%%%%%%%%%%%%%%%%%%%%%%%%%%%%%%%%%%%
%%%%%%%%%%%%%%%%%%%%%%%%%%%%%%%%%%%%%%%%%%%%%%%%%%%%%%%%%%%
%%%%%%%%%%%%%%%%%%%%%%%%%%%%%%%%%%%%%%%%%%%%%%%%%%%%%%%%%%%

\section{Einsatztaktische Überlegungen bei verunfallten Patienten}


\TextSlide{Unfallort}{%
Bevor die Patientenversorgung beginnen kann, müssen am
Unfallort einige Vorkehrungen für den weiteren Einsatzablauf
getroffen werden. Die wichtigsten Aufgaben sind:
\begin{itemize}
  \item \EE{Absichern} der Unfallstelle und Selbstschutz
  (Blaulicht, Warnblinkanlage, Uniformjacke anziehen, evtl.
  Helm aufsetzen) 
  \item Evtl. Nachalarmierung weiterer
  Einsatzkräfte: weitere Rettungsmittel bei mehreren
  Verletzten, Feuerwehr, Polizei (Alarmierung ausschließlich
  über das Journal, Meldung der Anzahl an Verletzten) 
  \item
  Bei mehreren Verletzten Triage
  (\FullPrettyRef{topic:triage})
  \item
  Eingeklemmte Personen werden im Fahrzeug stabilisiert bis
  die Feuerwehr die Verunfallten rettet; Helme tragen! 
  \item
  Crash-Rettung nur bei Fremdgefährdung, \Abk{HWS} stabilisieren!
\end{itemize}
}{
\begin{itemize}
  \item Absichern
  \item Evt. Nachalarmierung
  \item Triage
  \item Stabilisierung bei eingeklemmten Personen
\end{itemize}
}{}{}{}


\TextSlide{Unfallmechanismus}{%
Eine rasche, überblicksmäßige Beurteilung des
Unfallgeschehens erfolgt bereits im Szenenüberblick bzw.
indirekt in der Ersteinschätzung des Patienten. Der
 Unfallmechanismus kann dabei wichtige Hinweise
auf mögliche bzw. wahrscheinliche Verletzungen des Patienten
geben. Eine detaillierte Analyse des Unfallhergangs kann
durch die
Anamnese erfolgen. Zusätzlich kann das Unfallszenario
fotographiert werden. Primär müssen folgende Fragen geklärt
werden:
\begin{itemize}
    \item \emph{Was ist passiert?}
    \item \emph{Wie wurde der Patient verletzt?}
\end{itemize}
}{
\begin{itemize}
  \item Was?
  \item Wie?
\end{itemize}
}{}{}{}

\Fazit{Es ist sinnvoll, nach Möglichkeit den Unfallmechanismus 
mittels Fotos zu dokumentieren, um die Bilder der weiterbetreuenden 
Stelle zeigen zu können. Dies kann helfen, um bei der klinischen 
Versorgung des Patienten Prioritäten besser setzen zu können.}

\subsection{Besondere Unfallmechanismen}

\TextSlide{Übersicht}{
Im Folgenden werden einige Kennzeichen der wichtigsten
Unfallmechanismen vorgestellt.
\begin{itemize}
  \item Verkehrsunfälle
  \item Unfälle mit Fahrzeugen
  \item Frontalzusammenstöße
  \item Seitenaufprall
  \item Fußgänger vs. Fahrzeug
  \item Sturz aus großer Höhe
  \item Stich- und Schussverletzungen
\end{itemize}
}{
\begin{itemize}
  \item Verkehrsunfälle
  \item Unfälle mit Fahrzeugen
  \item Frontalzusammenstöße
  \item Seitenaufprall
  \item Fußgänger vs. Fahrzeug
  \item Sturz aus großer Höhe
  \item Stich- und Schussverletzungen
\end{itemize}
}{}{}{}

\TextSlide{Verkehrsunfälle}
{ Zur Beurteilung von Verletzungen durch Verkehrsunfälle
werden in der Literatur drei Faktoren genannt:
\begin{itemize}
  \item Ausmaß der Verformung des Fahrzeugs (Hinweis auf die
  beteiligten Kräfte)
  \item Beschädigung von Teilen in der
  Fahrzeugkabine (Hinweis für Aufschlagpunkt des Körpers im
  Fahrzeug) 
  \item Verletzungsmuster des
  Patienten (Hinweis darauf,
  welche Körperteile direkt aufgeprallt sind)
\end{itemize}
Es gilt weiters heraus zu finden, ob der Patient angegurtet
war und ob es Airbags gibt, die beim Aufprall ordnungsgemäß
ausgelöst wurden. 
\cite{ITLSde6,RedelsteinerHRNS1}
}{
\begin{itemize}
  \item Fahrzeug
  \item Teile im Fahrzeug
  \item Patient
  \begin{itemize}
    \item Verletzungsmuster
    \item Patient angegurtet?
    \item Airbags ausgelöst?
  \end{itemize}
\end{itemize}
}{}{}{}



\TextSlide{Unfälle mit Fahrzeugen}{%
Für Unfälle mit Fahrzeugen gilt allgemein, dass die
Geschwindigkeit des Fahrzeugs vor dem Unfall eine große
Auswirkung auf den Schweregrad des Unfalls und das Ausmaß der
entstandenen Verletzungen hat.
Je schneller das Fahrzeug unterwegs war, desto schwerer kann
die Verletzung sein. \Z{Dabei steigt die Bewegungsenergie
des Fahrzeugs  mit dem
Quadrat der Geschwindigkeit an. Eine Verdoppelung der
Geschwindigkeit  hat daher eine Vervierfachung der Energie
zur Folge, welche im Falle eines Unfalls frei wird und die
Verformung verursacht.}
Im Falle eines Unfalls (Aufprall auf ein Hindernis)
überträgt sich die Richtung der Gewalteinwirkung direkt auf
den Patienten. Ein Frontalzusammenstoß führt somit dazu,
dass der Patient von vorne verletzt wird, ein Seitenaufprall
führt zu Verletzungen auf der jeweiligen Körperseite des
Patienten. Von besonders schweren Verletzungen (SHT,
Wirbelsäulentrauma, Thoraxtrauma, Bauchtrauma, \ldots) ist
auszugehen, wenn der Patient aus dem Fahrzeug geschleudert wird.
Das Riskio für derart schwere Verletzungen ist um 300\PC*
erhöht \cite{BoehmerTaschRett6}.

Gibt es bei einem Unfall mehrere Verletzte und ist ein
Patient schwer verletzt, eingeklemmt oder sogar getötet,
muss man davon ausgehen, dass auch die anderen Unfallopfer
schwer verletzt sind. Diese Annahme gilt, bis bei
weiteren Untersuchungen (im Spital) das Gegenteil
bewiesen wurde! Stellt sich also bei der Behandlung eines
Patienten eine schwere Verletzung heraus, muss es zu einer
neuen Lageeinschätzung kommen.

Bei der Beurteilung des Verletzungsmusters muss bedacht
werden, dass es durch einen Zusammenstoß aufgrund der
Kraftübertragung zu einer ganzen Kette an Stößen kommt:
\begin{itemize}
  \item Kraftübertragung beim Aufprall des Fahrzeugs
  \item Kraftübertragung beim Aufprall des Körpers
  \item Kraftübertragung beim Aufprall der Organe innerhalb
  des Körpers.
  
  Beispiel für die Kraftübertragung innerhalb des Körpers:
  Prallt der Patient mit dem Knie gegen das Lenkrad oder das
  Amaturenbrett, so wird die Kraft (in Richtung der
  Gewalteinwirkung) über den Oberschenkel auch auf die Hüfte
  übertragen. Bei Knieverletzungen durch den Aufprall in
  einem Fahrzeug muss daher auch an mögliche Oberschenkel-
  und Beckenfrakturen gedacht werden.
\end{itemize}
Weiters kann es zu Stößen durch im Fahrzeug herumfliegende
Teile kommen. 
\cite{ITLSde6,RedelsteinerHRNS1,BoehmerTaschRett6} 
}{
\begin{itemize}
  \item Geschwindigkeit: Je schneller das Fahrzeug unterwegs
   war, desto schwerer die Verletzungen.
  \item Aus dem Fahrzeug geschleudert: Risiko  schwerer Verletzungen ↑↑↑
  \item Kraftübertragung:
  \begin{itemize}
    \item Aufprall des Fahrzeugs
    \item Aufprall des Körpers
    \item Aufprall der Organe
  \end{itemize}
\end{itemize}
}{}{}{}


\clearpage
\TextSlide{Frontal- und Heckzusammenstöße}{%
Bei Frontalzusammenstößen wird die Kraft von vorne auf das
Unfallopfer übertragen. Der Insasse wird aufgrund der
Trägheit nach vorne geschleudert und schlägt daher u.\,U.
auf das Lenkrad oder die Frontscheibe auf. Im Allgemeinen
gilt: War der Patient nicht angegurtet, ist mit schweren
Verletzungen zu rechnen! Aufgrund der Kraftrichtung von
vorne sind folgende schwere Verletzungen typisch:
Thoraxtrauma mit Herzkontusion und Pneumothorax,
\Abk{HWS}-Trauma, Leber- und Milzruptur, Beckenfraktur. 

Ein Heckaufprall hat für den
Patienten meistens die gleichen Folgen wie ein
Frontalzusammenstoß. In diesem Fall wird der Patient nicht
aufgrund der Trägheit nach vorne geschleudert, sondern
aufgrund der von hinten wirkenden Kraft. Beim Heckaufprall
kann es daher ebenso zu HWS-Trauma (Schleudertrauma) oder
Thoraxtrauma (Brust wird gegen Lenkrad gedrückt) kommen. Die
von hinten auf den Patient einwirkende Kraft wird zu einem
Großteil vom Sitz abgefedert.

Es gilt in beiden Fällen besonders zu beachten:
\begin{itemize}
    \item Ein Knieanstoß am Lenkrad oder Armaturenbrett
    kann in der Folge zu Oberschenkel- und Beckenfrakturen
    führen.
    \item Eine eingedrückte, leicht nach außen gewölbte
    Windschutzscheibe  (evtl. Spinnennetzmuster) deutet auf
    einen Aufprall des Kopfes gegen die Scheibe hin. SHT und
    HWS-Trauma sind zu befürchten.
    \item Sichtbare Veränderungen des Radstandes (z.\,B. Rad
    zur Seite geknickt, Achsenverschiebung um mehr als
    30\CM* am Unfallfahrzeug) spiegeln ein hohes Ausmaß an
    Verformung wieder \cite{BoehmerTaschRett6} . Es ist
    daher generell mit schweren Verletzungen zu rechnen.
    \item Faustregel: Fahrzeugverformung um mehr als 50\CM*
    sind als starke Deformierung zu werten.
    \cite{BoehmerTaschRett6} 
\end{itemize}
}{
\begin{itemize}
  \item Kraft von vorne bei
  \begin{itemize}
    \item Frontalzusammenstoß und
    \item Heckaufprall
  \end{itemize}
  \item Beachte:
  \begin{itemize}
  	\item Knieanstoß
  	\item Windschutzscheibe
  	\item Radstand
  	\item Fahrzeugverformung
  \end{itemize}
\end{itemize}
}{}{}{}


\TextSlide{Seitenaufprall}{%
Beim Seitenaufprall wirken die Kräfte seitlich auf den
Patienten ein. Es kommt daher häufig zu
HWS-Verrenkungen, Thoraxtrauma, Verletzungen der Aorta,
seitenabhängig Leber- oder Milzrupturen und/oder Frakturen
des Beckens \cite{BoehmerTaschRett6}. Wurden im Fahrzeug
Seiten- und Kopfairbags ausgelöst, ist zu beachten ob der
Patient in diese (oder in die entgegengesetzte) Richtung
geschleudert wurde. 
}{
\begin{itemize}
  \item Kraft von der Seite
  \item Beachte Kopf- und Seitenairbags
\end{itemize}
}{}{}{}



\Text{Fußgänger vs. Fahrzeug}{%
Stellt man im Szeneüberblick fest, dass
 ein \EE{Fußgänger von einem Fahrzeug (frontal) erfasst}
 wurde, ist zunächst immer von schweren Verletzungen
 (insbesondere SHT und Wirbelsäulentrauma) auszugehen.
 In \cite{BoehmerTaschRett6} sind folgende
 Faustregeln angegeben:
\begin{itemize}
  \item Bis 30\,\MetricUnit{\nicefrac{km}{h}}:
  Lebensbedrohliche Verletzungen möglich\\
  (Eine
  Fahrzeuggeschwindigkeit
  von mehr als 30\,\MetricUnit{\nicefrac{km}{h}} kann bei trockener
  Fahrbahn an Bremsspuren, die länger als
  \VonBis{10}{20}\Meter* sind, erkannt werden.)
  \item \Z{Bis 50\,\MetricUnit{\nicefrac{km}{h}}:
  Aufprall des Kopfes auf die Kühlerhaube}
  \item \Z{Bis 70\,\MetricUnit{\nicefrac{km}{h}}: Aufprall
  des Kopfes auf die Windschutzscheibe}
  \item \Z{Mehr als 70\,\MetricUnit{\nicefrac{km}{h}}:
  Fußgänger wird über das Fahrzeugdach geworfen (gilt für
  PKW)}
\end{itemize}
}{
\SymbolText
}{}{}{}

% \begin{table}[b]
% \begin{WidePar}
% \Captionof{table}{Bilderserie: Unfallmechanismen}
% \begin{tabularx}{\linewidth}{XXXX}
% \includegraphics[width=\linewidth]{./images/teaser/doa.jpg}
% 
% \Captionof{figure}{Oft der Verlierer: Der
% Fußgänger. \SourcePicture{Hauer}}
% &
% &
% &
% \\\bottomrule
% \end{tabularx}
% \end{WidePar}
% \end{table}

\begin{PictureSeriesWide}{}{Bilderserie: Unfallmechanismen}
\PictureSeriesRow{3}{./images/teaser/doa.jpg}
{Oft der Verlierer: Der
Fußgänger. \SourcePicture{Hauer}}
{empty}{}
{empty}{}
{empty}{}
\end{PictureSeriesWide}



\TextSlide{Sturz aus großer Höhe}{%
Bei Stürzen aus großen
Höhen (z.\,B. Leiter, Gerüst, Balkon) ist im Szeneüberblick
an folgende Fragestellungen zu denken:
\begin{itemize}
  \item Aus welcher Höhe ist der Patient gestürzt?
  
  Ab einer Fallhöhe von 3\Meter* ist mit schweren
  Verletzungen zu rechnen, ab einer Fallhöhe von
  \E{6\Meter*} ist von einem \E{Polytrauma}
  (\FullPrettyRef{topic:polytrauma})
auszugehen \cite{BoehmerTaschRett6}.
  \item Mit welchem Körperteil ist er auf dem Boden
  aufgeprallt?
  
  Generell ist wieder mit SHT und Wirbelsäulentrauma zu
  rechnen. Bei einem Aufprall auf die gestreckten Beine ist
  zusätzlich an Frakturen der unteren Extremitäten (z.\,B.
  Fraktur der Fersenbeine) sowie an Stauchungen der
  Wirbelsäule zu denken. Kinder schlagen oft mit dem Kopf
  auf, weshalb hier SHT häufig anzutreffen sind.
  \item Wie ist der Boden beschaffen (Beton, Wiese etc.)?
  Ist der Patient auf herumliegende Gegenstände gefallen?
\end{itemize}
}{
\begin{itemize}
  \item Beachte:
  \begin{itemize}
    \item Fallhöhe
    \item zuerst aufprallender Körperteil
    \item Oberflächenbeschaffenheit des Bodens
  \end{itemize}
  \item Rechne mit
  \Abk{HWS}-Trauma, Wirbelsäulentrauma, SHT, Beckenfraktur
\end{itemize}
}{}{}{}


\TextSlide{Stich- und Schussverletzungen}{%
Stich- bzw. Pfählungsverletzungen und Schussverletzungen
werden zu \EE{penetrierende Verletzungen} zusammen
gefasst.\footnote{\Z{lat. penetrare: eindringen,
durchdringen.}}

Besonderheiten von Stichverletzungen sind, dass sie meistens
unspektakulär aussehen. Sie sind oft klein und nur schwach
blutend. Dennoch können sie im Körperinneren großen Schaden
anrichten. Um das Verletzungsausmaß von Stichwunden richtig
beurteilen zu können, müssen folgende Aspekte beachtet
werden:
\begin{itemize}
  \item Welche Körperteile sind betroffen? (Welche Organe
  liegen direkt unter der Einstichstelle?)
  \item Gibt es mehrere Einstichstellen? Wie viele?
  \item Wie lang war die Klinge?
  \item Wie war der Einstichwinkel? (Welche Organe könnten
  noch betroffen sein?)
  
  Beispiel: Ein schräg in den Oberbauch eindringender
  Gegenstand kann auch Organe im Thorax verletzen!
  \item Wurde nach dem Einstich die Lage des Gegenstands
  im Körper verändert?
\end{itemize}

Die Besonderheit von Schussverletzungen ist, dass
das Geschoss mit sehr großer Geschwindigkeit in den Körper
eindringt. Die dadurch entstehende Druckwelle verursacht
schwere sekundäre Verletzungen. Schussverletzungen sind in
Österreich verhältnismäßig selten.

Im Körper steckende Messer, Stichwaffen und andere
\EE{Fremdkörper dürfen niemals präklinisch entfernt werden}!
Eventuell verhindert ein noch steckender Fremdkörper gerade
noch eine unkontrollierte Blutung. Eine Entfernung darf
daher erst unter kontrollierten Bedingungen \EE{in einem
Operationssaal} erfolgen. Für den Transport muss man den
\EE{Fremdkörper mittels eines Verbandes sichern}, um ein
Verrutschen und Folgeverletzungen zu verhindern.
}{
\begin{itemize}
  \item Können harmlos aussehen, aber trotzdem
  lebensbedrohlich sein!
  \item Fremdkörper nicht entfernen!
  \item Fremdkörper mittels Verband sichern!
\end{itemize}
}
{./images/teaser/blut-tatort-hinterhof.jpg}
{}
{}

\begin{PictureSeriesWide}{}{Bilderserie: \EE{Fremdkörper}}
\PictureSeriesRowWide{3}
{./images/trauma/fremdkoerper-stich-fixierung.jpg}
{Fixierung
eines penetrierenden Fremdkörpers. \SourcePicture{Hauer}}
{./images/trauma/fremdkoerper-stich-messer-01.jpg}
{Stich mit einem Küchenmesser. 
Das Messer verbleibt im Körper {\ldots}~
\SourcePicture{Gabriel}}
{./images/trauma/fremdkoerper-stich-messer-02.jpg}
{{\ldots} bis in den Operationssaal.
\SourcePicture{Gabriel}}
{}
{}
\end{PictureSeriesWide}


\subsection{Zusammenfassung}
Aus allen einsatztaktischen Überlegungen ergibt sich ein
Bündel an Erstmaßnahmen, welche bei Unfällen durchgeführt
werden müssen.


% M %%%%%%%%%%%%%%%%%%%%%%%%%%%%%%%%%%%%%%%%%%%%%%%%%%% M %
% Unfälle
\ProcedurePrintDefault{TY14101C}



%%%%%%%%%%%%%%%%%%%%%%%%%%%%%%%%%%%%%%%%%%%%%%%%%%%%%%%%%%%
%%%%%%%%%%%%%%%%%%%%%%%%%%%%%%%%%%%%%%%%%%%%%%%%%%%%%%%%%%%
%%%%%%%%%%%%%%%%%%%%%%%%%%%%%%%%%%%%%%%%%%%%%%%%%%%%%%%%%%%
\clearpage % TEMP temp clearpage
\section{Verletzungen der Knochen: Frakturen}

%%%%%%%%%%%%%%%%%%%%%%%%%%%%%%%%%%%%%%%%%%%%%%%%%%%%%%%%%%%
%%%%%%%%%%%%%%%%%%%%%%%%%%%%%%%%%%%%%%%%%%%%%%%%%%%%%%%%%%%
%%%%%%%%%%%%%%%%%%%%%%%%%%%%%%%%%%%%%%%%%%%%%%%%%%%%%%%%%%%
\label{topic:frakturen}%

\subsection{Allgemeines}

%\paragraph{Arten}
\TextSlide{{\TxBeschreibung} und Arten}
{%
Ein \DefinitionInPara{\T{Knochenbruch} 
(\Lang{Syn.} \T[Fraktur]{Fraktur}\Z{, \Lang{lat.} \IX{Fractura} (\I{Fract.})}) 
ist eine Kontinuitätsunterbrechung eines Knochens 
mit Bildung von Fragmenten (Bruchstücken) \cite{Pschy259}.}
%
Grundsätzlich wird zwischen zwei Arten von Knochenbrüchen
unterschieden: Beim \EE{geschlossenen Knochenbruch} ist die
Haut über der Bruchstelle noch intakt, während beim
\EE{offenen Knochenbruch}\ZFootnote{Offener Bruch: \I{Fractura aperta} (\I[Fract.!apert.]{Fract. apert.})} 
aufgrund der Gewalteinwirkung von
außen die Haut über der Bruchstelle verletzt ist. Der Knochen muss dabei nicht
herausragen.\footnote{\Z{Die Literatur ist bei der
Definition offener Knochenbrüche nicht einheitlich. In
manchen Büchern, z.\,B. in \cite{Gorgass7}, wird ein offener
Knochenbruch durch eine direkte Verbindung zwischen Wunde
und Bruchstelle
beschrieben, in anderen Büchern genügt es, wenn über
der Bruchstelle eine Hautverletzung vorliegt. Wieder andere
Bücher, z.\,B. \cite{Strauss0000}, geben drei Schweregrade
für offene Frakturen an. Die genau Definition spielt für die
Praxis keine Rolle. Das Achten auf Keimfreiheit ist in
jedem der beschriebenen Fälle unbedingt notwenidg!}}}{%
\begin{itemize}
    \item Geschlossener Knochenbruch
    \item Offener Knochenbruch
\end{itemize}
}{}{}{}

\TextSlide{{\TxSymptome}: Frakturzeichen}
{\label{topic:frakturzeichen}%
Erkennbar ist ein Knochenbruch anhand von klassischen Frakturzeichen. 
Dabei gibt es sichere und unsichere Frakturzeichen, welche in
\FullPrettyRef{tab:frakturzeichen} angeführt sind. Nicht immer kann 
die Diagnose eines Bruches gestellt oder ausgeschlossen werden.
}
{
\begin{itemize}
  \item Sichere Frakturzeichen
  \item Unsichere Frakturzeichen
  \item \FullPrettyRef{tab:frakturzeichen}
\end{itemize}
}
{}
{}

\Layout{47}{}
{}
{}
{\index{Frakturzeichen}%
\begin{tabularx}{\linewidth}[H]{XX}\addlinespace\toprule\addlinespace
\multicolumn{1}{c}{\TabHeading{Sichere Frakturzeichen}}
 &
\multicolumn{1}{c}{\TabHeading{Unsichere Frakturzeichen}}
\\\addlinespace\midrule
\begin{itemize}
    \item Abnorme Beweglichkeit
    \item Fehlstellung
    (abnorme Stellung, Verkürzung, Treppenbildung)
    \item Knochenreiben (Krepitus) \footnote{\so{NICHT}
    absichtlich ausprobieren!}
    \item Sichtbare freie Knochenteile (= offener Bruch)
\end{itemize}
&
\begin{itemize}
    \item Schmerz
    \item Schwellung
    \item Hämatom
    \item Gestörte Funktion / Bewegungsunfähigkeit
\end{itemize}
\\\bottomrule
\end{tabularx}
}
{Frakturzeichen\label{tab:frakturzeichen}}
{}


\TextSlide{\TxGefahren}{%
\begin{itemize}
  \item \Ca{Starke Blutungen} bis hin zum Schock 
  Je nach
  betroffenen Knochen kann es zu unterschiedlich großen
  Blutverlusten kommen:
  \begin{labeling}[~~:]{Unterschenkel}
    \item[Oberarm] \VonBis{100}{\EE{800}}\Ml*
    \item[Unterarm] \VonBis{50}{\EE{400}}\Ml*
    %	\item[Rippen: \EE{2000-3000\Ml*
   \item[\Ca{Becken}] \EE{\VonBis{500}{\Ca{5\,000}}}\Ml*
   \item[\Ca{Oberschenkel}] \EE{\VonBis{300}{\Ca{2\,000}}}\Ml* (Schaftfraktur!)
   \item[Unterschenkel] \VonBis{100}{\EE{1\,000}}\Ml*
  \end{labeling}
  Die Blutung kann aus dem Knochenmark, zerrissener Muskulatur oder aus einem
  verletzten Gefäß kommen.
  \item \Ca{Fettembolie} (v.\,a. bei Fraktur langer Röhrenknochen): Lungenembolie durch ausgeschwemmtes Fett
  \item Begleitende \EE{Nervenverletzungen}
  \item \EE{Infektionsgefahr}: Bei offenen Frakturen ist
  Keimfreiheit be\-sonders wichtig, da  die Gefahr einer
  lebenslangen Knochen\-marks\-ent\-zündung
  \Z{(Osteomyelitis)} besteht!
  Wenn mögllich soll eine Wundfolie (z.\,B.
  \index{OpSite}\T{OpSite}\texttrademark ) verwendet werden.
  \item Kombination von mehreren Frakturen $\rightarrow$
  großer Blutverlust 
  \item Schwere Muskelquetschung kann zu
  einem Sauerstoffmangel im Muskel führen
  (Kompartmentsyndrom)
\end{itemize}
}{%
\begin{itemize}
  \item \Ca{Starke Blutung}
  \item \EE{Fettembolie}
  \item Nervenverletzungen
  \item Infektion
  \item Kombination mehrerer Frakturen
  \item Muskelquetschungen
\end{itemize}
}{}{}{}

Auch bei einer harmlos anmutenden Fraktur muss ein
gründlicher Traumacheck durchgeführt werden, damit keine
Begleitverletzungen übersehen werden. Dafür ist es unbedingt
notwendig, dass der Patient entkleidet
wird.\footnote{Nachdem der Patient für die Untersuchung
entkleidet wurde, muss er nach der Untersuchung zugedeckt
werden (Wärmeerhaltung)!} Beim Entfernen der Schuhe ist
darauf zu achten, dass eine Person das Bein fixiert
während eine andere Person den Schuh öffnet und schließlich 
mittels Stiefelgriff (siehe
\prettyref{topic:vacuumbeinschiene-anwendung}) vorsichtig
vom Fuß abzieht.

%\paragraph{Frakturzeichen} 



\begin{PictureSeriesWide}{}{Bilderserie: \SlideHeading{Frakturen}}
\PictureSeriesRowWide{3}
{./images/fraktur_fehlstellung_nativ.jpg}
{Fehlstellung einer
Unterschenkelfraktur nach einem Verkehrsunfall (PKW gegen
Fußgänger) in natura {\ldots} \SourcePicture{Hauer}}
{./images/fraktur_fehlstellung_xray.jpg}
{{\ldots} und in Röntgendarstellung \SourcePicture{Hauer}}
{./images/trauma/fraktur-offen-finger.jpg}
{Nicht immer sind offene Frakturen
spektakulär: Offene Fraktur des Mittelfingers
\SourcePicture{Ben Stephenson, Cleveland, OH, USA.
Lizenz: CC-BY 2.0}}
{}
{}
\end{PictureSeriesWide}


% \begin{PictureSeriesWide}{}{Bilderserie: \EE{Frakturen}}
% \PictureSeriesRowWide{3}
% {}
% {}
% {./images/trauma/fraktur-us-nativ-rotated.jpg}
% {Unterschenkelfraktur
% nach Verkehrsunfall 
% \SourcePicture{Hauer}}
% {./images/trauma/fraktur-us-roe.jpg}
% {Röntgenbild
% der Unterschenkelfraktur \SourcePicture{Hauer}}
% {}
% {}
% \end{PictureSeriesWide}

% M %%%%%%%%%%%%%%%%%%%%%%%%%%%%%%%%%%%%%%%%%%%%%%%%%%% M %
% Fraktur
\ProcedurePrintDefault{TT14200C}

\subsection{Spezielle Knochenbrüche}

\TextSlide{\TxBeschreibung}{
Die beiden im Rettungsdienst am häufigsten anzutreffenden,
speziellen Knochenbrüche sind die Schlüsselbeinfraktur und
die Schenkelhalsfraktur.
}{
\begin{itemize}
  \item Schlüsselbeinfraktur
  \item Schenkelhalsfraktur
\end{itemize}
}{}{}{}

\paragraph{Schlüsselbeinfraktur} Die Schlüsselbeinfraktur
entsteht meistens, wenn der Patient mit ausgestreckten Armen
stürzt. Die große Gefahr dieser Fraktur liegt darin, dass
die Schlüsselbeinarterie und -vene verletzt werden können.
Die Ruhigstellung erfolgt einerseits durch ein Armtragetuch
(Dreiecktuch) und andererseits durch Fixation des Arms am
Oberkörper durch ein zweites Dreiecktuch.

\paragraph{Schenkelhalsfraktur} Bei der Schenkelhalsfraktur
handelt es sich um einen Bruch des Oberschenkelhalses. Meist
sind ältere Menschen betroffen, die auf die Hüfte gestürzt
sind. Symptome sind:
\begin{itemize}
	\item Starker Bewegungs- oder Stauchungsschmerz
	\item Betroffenes Bein ist nach außen gedreht und verkürzt
	\item Unfallhergang meist Sturz (auch an Kollaps denken!)
\end{itemize}

Die Versorgung erfolgt durch Rettung mittels Schaufeltrage
und Ruhigstellung mittels Vakuummatratze. An große
Blutverluste denken (Kontrolle der Vitalparameter)!
Eventuell kann eine Schmerztherapie (Notarzt) notwendig
sein.

%%%%%%%%%%%%%%%%%%%%%%%%%%%%%%%%%%%%%%%%%%%%%%%%%%%%%%%%%%%
%%%%%%%%%%%%%%%%%%%%%%%%%%%%%%%%%%%%%%%%%%%%%%%%%%%%%%%%%%%
%%%%%%%%%%%%%%%%%%%%%%%%%%%%%%%%%%%%%%%%%%%%%%%%%%%%%%%%%%%
\section{Verletzungen der Gelenke: Verrenkung und Verstauchung}


\subsection{Verstauchung, Bänderzerrung und Bänderriss}

\label{topic:verstauchung-tra}

\Text{\TxBeschreibung}
{%
\DefinitionInPara{Eine \T{Verstauchung} 
(\Lang{Syn.} \T{Distorsion}, \Z{\Lang{lat.} \I{Distorsio} (\I{Dist.})}) 
ist eine 
Zerrung bzw. Überdehnung des Kapsel-/ Bandapparates eines Gelenks.}
%
Bei der Verstauchung (\T{Distorsion}) handelt es sich um
eine Gelenksverletzung, bei welcher ein Knochen durch 
stumpfe Gewalteinwirkung kurz verschoben oder verdreht wird.
Im Gegensatz zur Verrenkung springt der Knochen aber von
selbst wieder in seine ursprüngliche Stellung im Gelenk
zurück.

Da Gelenke durch Bänder vor unnatürlichen bzw. übermäßigen
Bewegungen geschützt werden, können diese sowie die
Gelenkskapsel bei einer Distorsion in Mitleidenschaft
gezogen werden und bei ausreichender Gewalteinwirkung reißen.
}



\TextSlide{\TxSymptome}{%
Bei einer Distorsion klagt der Patient über \E{Schmerzen} im
Gelenksbereich. Je nach Unfallschwere kann der Patient das
entsprechende Gelenk evtl. \E{eingeschränkt bewegen}. Durch
die Weichteilverletzung kommt es zu einem \EE{Anschwellen}
des umliegenden Gewebes und zu einem \E{Bluterguss}
(Hämatom). Falls es zu einem Bänderriss gekommen ist,
berichten der Patient oder die Unfallzeugen meist über ein
hörbares \E{Schnalzen}.

\bigskip\bigskip

}{
\begin{itemize}
  \item Schmerzen, Schwellung, Bluterguß (Hämatom)
  \item Evtl. Bewegungseinschränkungen
  \item Hörbares Schnalzen beim Bänderriss
\end{itemize}
}{}{}{}


% M %%%%%%%%%%%%%%%%%%%%%%%%%%%%%%%%%%%%%%%%%%%%%%%%%%% M %
% Verstauchung
\ProcedurePrintDefault{TT14031C}





\subsection{Verrenkung}
\label{topic:verrenkung}

\Text{\TxBeschreibung}
{%
\DefinitionInPara{Eine \T{Verrenkung} 
(\Lang{Syn.} \T{Luxation}, \Z{\Lang{lat.} \I{Luxatio} (\I{Lux.})})  
ist eine Gelenkverletzung,
bei der zwei das Gelenk bildende Knochenenden aus ihrer
normalen Stellung zueinander verschoben werden. Dabei kommt
es meist zu Begleitverletzungen an der Gelenkkapsel bzw. den
Gelenkbändern \cite{Gorgass7}.}
%
Die Verrenkung (\T{Luxation}) ist eine Gelenksverletzung,
bei der -- im Gegensatz zur Verstauchung -- der verschobene
oder verdrehte Knochen nicht von selbst in das Gelenk zurück
gesprungen ist. Die betroffene Extremität kann daher nicht
mehr bewegt werden. Bei einer Luxation dürfen keine
eigenständigen Einrenkversuche unternommen werden!

Durch das inadäquate Verhältnis zwischen Kopfgröße und
Pfannengröße im Schultergelenk ist dies das am häufigsten
luxierte Gelenk im menschlichen Körper. \Z{45\PC* aller
Luxationen entfallen auf das
Schultergelenk \cite{Gorgass7}.}
}


\TextSlide{\TxSymptome}{%
Der Patient klagt über Schmerzen im Bereich des Gelenks,
durch die Weichteilverletzung kommt es zu einer Schwellung
der umliegenden Gewebe und zu einem Hämatom. Klinisch kann
man eine abnorme Gelenksstellung, welche einen federnd
fixierten Charakter hat, feststellen. Der Patient kann das
entsprechende Gelenk nicht bewegen.

\bigskip

\bigskip
}{
\begin{itemize}
  \item Schmerzen, Schwellung, Bluterguss
  \item Abnorme Gelenkstellung (federnd fixiert), bewegungsunfähig
\end{itemize}
}{}{}{}


% M %%%%%%%%%%%%%%%%%%%%%%%%%%%%%%%%%%%%%%%%%%%%%%%%%%% M %
% Verrenkung
\ProcedurePrintDefault{TT14032C}



%%%%%%%%%%%%%%%%%%%%%%%%%%%%%%%%%%%%%%%%%%%%%%%%%%%%%%%%%%%
%%%%%%%%%%%%%%%%%%%%%%%%%%%%%%%%%%%%%%%%%%%%%%%%%%%%%%%%%%%
%%%%%%%%%%%%%%%%%%%%%%%%%%%%%%%%%%%%%%%%%%%%%%%%%%%%%%%%%%%
\clearpage % TEMP temp clearpage
\section{Schädel-Hirn-Trauma (SHT)}
%Folge von Gewalteinwirkung auf den  Schädel
\label{topic:sht}%



\Layout{60}{\TxBeschreibung{} und Arten}
{%
\DefinitionInPara{Ein \T[Schädel-Hirn-Trauma|see{SHT}]{Schädel-Hirn-Trauma} 
(\Lang{Abkz.} \T{SHT}) ist eine
Gewalteinwirkung auf den Kopf, welche 
Funktionsstörungen und Verletzungen des Gehirns hervorruft \cite{Gorgass7}.}
%
Die häufigsten Ursachen von Schädel-Hirn-Traumata (SHT) sind
Verkehrsunfälle, Stürze und Gewaltverbrechen. Je nach
Schwere und Mechanismus können verschiedene  Formen des \Abk{SHT}
auftreten. Man unterscheidet das \EE{gedeckte \Abk{SHT}}, bei dem
die harte Hirnhaut (Dura mater) nicht eröffnet wird, sowie
das \EE{offene \Abk{SHT}}, bei dem die harte Hirnhaut eröffnet
wird. Bei beiden Arten kommt es zu Bewusstseinsstörungen.
}{
\begin{compactitem}
  \item Leitsymptom: \EE{Bewusstseinsstörung}.
  \item Arten:
  \begin{compactitem}
    \item Gedecktes \Abk{SHT}: ohne  Eröffnung der harten Hirnhaut
    \item Offenes \Abk{SHT}: mit  Eröffnung der harten Hirnhaut
  \end{compactitem}
\end{compactitem}

}
{./images/teaser/imgp0722_00500.jpg}
{}
{GABS.}

\Fazit{Das Leitsymptom des SHT
ist die \EE{Bewusstseinsstörung}.}

\TextSlide{Begleitende Verletzungen}
{%
\begin{itemize}
  \item \Abk{HWS}-Verletzung
  \item Verletzung der Kopfschwarte  (Rissquetschwunde)
  \item Frakturen von
  \begin{itemize}
    \item Schädeldach
    \item Schädelbasis (Blutung bzw. Liquoraustritt
    aus Ohren oder Nase, \I[Monokelhämatom]{Monokel-}
    oder \I{Brillenhämatom})
    \item Gesichtsschädel
    \item Unterkiefer
  \end{itemize}
\end{itemize}
}{
\begin{itemize}
  \item \Abk{HWS}-Verletzung
  \item Rissquetschwunden
  \item Frakturen
\end{itemize}
}{}{}{}




\Cave{Die Schwere eines \Abk{SHT} lässt sich nur bedingt von
den äußerlich sichtbaren  Verletzungen ableiten! Oft werden
im Krankenhaus Frakturen festgestellt,  ohne dass äußere
Wunden vorliegen. Eine \EE{sorgfältige Untersuchung} ist
daher bei jedem \Abk{SHT} Pflicht!}

\TextSlide{\TxGefahren}
{%
\begin{itemize}
  \item \EE{\Ca{Luzides Intervall} (lichtes Intervall)}:
  Die
  Symptome können eventuell erst Stunden nach der Verletzung
  einsetzen. Dazwischen ist der Patient unauffällig.
  Diese Zeitspanne wird \T[Intervall!luzides]{Luzides
  Intervall} genannt und ist sehr gefährlich, da man den
  Patienten leicht falsch einschätzen kann.
  \item \EE{Hirnblutung}:
  \index{Insult!blutiger}%
  \index{Hirnblutung}%
  \index{Blutung!intrakranielle}%
  Eine Hirnblutung (intrakranielle Blutung) kann sowohl in
  Folge eines \Abk{SHT} als auch spontan (z.\,B. durch Platzen eines
  Aneurysmas, \enquote{blutiger Insult}, \prettyref{topic:insult})
  auftreten.
  
  Aufgrund der inneren Blutung steigt der Druck im Gehirn
  an, es treten \II{Hirndruckzeichen} (pathologische
  Pupillenreaktion, Druckpuls). Im Extremfall kann es zu
  einer \I{Hirnstammeinklemmung} kommen. Patienten mit
  intrakranieller Blutung sind vital bedroht! Beim \Abk{SHT}
  ist zu beachten, dass die Symptome oft auch erst Stunden
  nach dem Unfall auftreten können (luzides Intervall).
  \item \EE{HWS-Verletzung}:
  Oft kommt es bei einer einem \Abk{SHT} zu einer
  Mitverletzung der Halswirbelsäule.
  \item \EE{Unregelmäßige Atmung}
\end{itemize}

}{%
\begin{itemize}
  \item \Ca{Luzides Intervall}
  \begin{itemize}
    \item Symptome erst Stunden nach Verletzung
  \end{itemize}
  \item Hirnblutung 
  \item \Abk{HWS}-Verletzung
  \item Unregelmäßige Atmung
\end{itemize}
}{}{}{}

\Layout{20}{Spezielle Diagnostik beim SHT}
{\label{topic:sht-diagnostik}%
Grundsätzlich ist beim Schädel-Hirn-Trauma die übliche
Trauma-Diagnostik durchzuführen. Die \E{neurologische
Beurteilung}, die Erhebung des \E{Unfallmechanismus} und die
suche nach \E{Begleit- oder Folgeverletzungen} sind
besonders zu betonen.

\begin{itemize}
  \item \EE{Unfallmechanismus}: Anhand des
  Unfallmechanismus muss ermittelt werden, ob es eine
  Verletzung der (Hals-) Wirbelsäule wahrscheinlich ist.
  Besonders bedenklich sind Sportverletzungen im Kopf- oder
  Nackenbereich, ein Trauma nach Sprung ins Wasser und
  Stauchungsverletzungen der \Abk{HWS} (z.\,B. Schlag auf
  Kopf;
  vgl.
  \FullPrettyRef{MK-m:einschaetzung-indikation-ws-immobilisation}).
  Ggfs. muss eine Immobilisation durchgeführt werden
  \item \EE{Neurologische Beurteilung}:
  Die neurologische Beurteilung ist beim \Abk{SHT} natürlich
  besonders wichtig und besonders genau durchzuführen. Ein
  \EE{Neurocheck} ist verpflichtend. Auffälligkeiten können
  Hinweise auf eine direkte Hirnschädigung oder eine Hirnblutung sein.
  
  Die
  \EE{Bewusstseinslage} ist besonders genau abzuklären: Eine
  Bewusstseinsstörunge kann auch andere
  Ursachen als den gegenständlichen Unfall haben (Blutzucker,
  Alkohol, Exsikkose, Schlaganfall, \ldots). Evtl. ist diese
  andere Ursache überhaupt erst der Auslöser für einen Unfall.
  
  Die \EE{Pupillen} können im Sinne von Hirndruckzeichen
  (\prettyref{topic:hirndruckzeichen}) pathologisch sein
  (ungleich weit, verlangsamte, einseitige oder fehlende
  Lichtreaktion).
  
  Daneben ist auch zu beurteilen, ob andere Hirndruckzeichen
  (Hypertonie mit Bradykardie, Strecksynergismen, \ldots)
  vorliegen, vgl. \FullPrettyRef{topic:hirndruckzeichen}.
  
  Weiters muss auf \EE{neurologische Ausfälle} (sensorisch,
  motorisch) in Folge einer \EE{\Abk{HWS}-Verletzung} geachtet
  werden.
  
  \Z{Die Schwere der Bewusstseinsstörung kann der erfahrene
  Helfer zusätzlich durch die Glasgow Coma Scale
  einschätzen.}
  \item \EE{Traumacheck} Ein Traumacheck ist natürlich
  verpflichtend durchzuführen
  (\FullPrettyRef{topic:traumacheck}).
  
  Besonders zu beachten ist die Gefahr einer hohen
  Querschnittslähmung, ebenso dass Schädel, \Abk{HWS} und
  \Abk{BWS} eine
  Linie bilden.
  
  \item \EE{Spitalsabklärung}: Jedes \Abk{SHT}, auch wenn die
  Diagnostik vor Ort unauffällig und der Patient gegenwärtig
  symptomfrei ist, ist im Spital abzuklären!
\end{itemize}


\Cave{\EE{Jedes \Abk{SHT} ist im Spital abzuklären!} Es besteht
die Gefahr dass sich der Patient gegenwärtig im
\protect\enquote{\EE{luziden Intervall}} 
befindet und sich sein Zustand erst in einigen Stunden, dann
aber schlagartig verschlechtert. Im Zweifelsfall wird der
Patient dort zur Beobachtung (24\Hour* lang) aufgenommen.}
}
{
\begin{itemize}
  \item Neurologische Beurteilung wichtig!
  \begin{itemize}
    \item Bewusstseinsstörungen können andere Ursachen
    haben!
    \item Hirndruckzeichen (Pupillen!)
  \end{itemize}
  \item Traumacheck
  \item Weitere Diagnostik im Spital, auch bei symptomfreien
  Patienten
  
  \Ca{Luzides Intervall}
\end{itemize}
}
{./images/trauma/pupillendifferenz.jpg}
{\Ca{Alarmzeichen:} Eine Pupillendifferenz kann ein
wichtiger Hinweis auf einen erhöhten Hirndruck sein!}
{}
{}


% M %%%%%%%%%%%%%%%%%%%%%%%%%%%%%%%%%%%%%%%%%%%%%%%%%%% M %
% SHT
\ProcedurePrintDefault{TS06XX0C}


\subsection{Offenes SHT}

\TextSlide{Hinweise}
{%
Hinweise auf ein offenes \Abk{SHT} sind:

\begin{itemize}
    \item Große RQW mit tast-/sichtbaren  Knochensplittern
    \item Hirnaustritt (schlechte Prognose)
    \item Blutungen aus Nase/Ohren
    \item Liquoraustritt, oft mit Blut vermischt
    ({\quotedblbase}Tupfertest{\textquotedblleft})
\end{itemize}
}{
\begin{itemize}
  \item Knochensplitter
  \item Hirnaustritt
  \item Blutungen
  \item Liquoraustritt
\end{itemize}
}{}{}{}


\subsection{Gedecktes SHT}
% EMHO: Taschenatlas RD spricht hier vom geschlossenen SHT
% (statt gedecktem SHT).

\TextSlide{\TxBeschreibung}{
Ein gedecktes (geschlossenes) Schädel-Hirn-Trauma ist ein
\Abk{SHT} ohne Eröffnung der harten Hirnhaut. Je nach
Schwere wird das \Abk{SHT} in drei Grade eingeteilt:

\begin{itemize}
\item \T{Gehirnerschütterung} {\dotfill} \Abk{SHT}
1{\textdegree}, Commotio cerebri
\item \T{Gehirnprellung} {\dotfill} \Abk{SHT}
2{\textdegree}, Contusio cerebri
\item \T{Gehirnquetschung} {\dotfill} \Abk{SHT}
3{\textdegree}
\end{itemize}
}{
\begin{itemize}
  \item Gehirnerschütterung {\hfill\ldots} \Abk{SHT}
  1{\textdegree}
  \item Gehirnprellung {\hfill\ldots} \Abk{SHT}
  2{\textdegree}
  \item Gehirnquetschung {\hfill\ldots} \Abk{SHT} 3{\textdegree}
\end{itemize}
}{}{}{}


\TextSlide{Grade des geschlossenen SHT}
{%
\subparagraph{Gehirnerschütterung -- SHT 1\textdegree}
\index{Gehirnerschütterung}

\Lang{Term.} \I{Commotio cerebri}

\noindent Die Gehirnerschütterung ist die leichteste Form des
Schädelhirntraumas. Leitsymptom ist ein Trauma mit
anschließender \EE{kurzer Bewusstlosigkeit}. Weiters kann es
zu flüchtigen neurologischen Ausfällen kommen, es kann auch eine
\I[Anmesie!retrograde]{retrograde Amnesie}
(Gedächtnisverlust über den Vorgang bzw. Unfall) bestehen.
Eventuell klagt der Patient über
Übelkeit erbrechen und Schwindel. Die Gehirnerschütterung
zieht keine bleibenden Schäden nach sich.

\subparagraph{Gehirnprellung -- SHT 2\textdegree}
\index{Gehirnprellung}

\Lang{Term.} \I{Contusio cerebri}

\noindent Die Gehirnprellung ist eine schwerere Form des
Schädel-Hirntraumas. Leitsymptom ist die
\EE{Bewusstlosigkeit länger
 als 1\Hour*}, beziehungsweise die Bewusstseinseintrübung
 <\,24\Hour*. Mittels weiterführender Diagnostik im Spital
 lassen sich die \EE{Hirnschädigungen} nachweisen. Es kann
 zu Nervenverletzungen kommen, dies kann zu vegetativen
 Störungen führen. Bleibende Schäden sind möglich.

\subparagraph{Gehirnquetschung -- SHT 3{\textdegree}}
\index{Gehirnquetschung}%
\label{topic:hirnblutung}%

\noindent Die Gehirnquetschung ist eine sehr schwere Form des
Schädelhirntraumas. Es kommt zu einer \EE{Bewusstlosigkeit
länger als 24\Hour*} und zu teils sehr ausgeprägten
Begleitsymptomen wie zum Beispiel Hirndruckzeichen
(\prettyref{topic:hirndruckzeichen}), Atem und
Kreislaufstörungen bis hin zum Atemstillstand durch
Einklemmen des Atemzentrums.

Es kann weiteres zu intrakraniellen Blutungen kommen
(epidural, subdural, subarachnoidal). Außerdem kommt es zu
verschiedenen neurologischen Ausfällen und Symptomen wie zum
Beispiel zerebrale Krampfanfälle. Mit \EE{bleibenden
Schäden} und Langzeitfolgen ist zu rechnen.

\bigskip

\bigskip

\SlideCompensation{1}
}{
\begin{itemize}
  \item \EE{Gehirnerschütterung -- 1\textdegree}
  \begin{compactitem}
    \item  Bewusstlosigkeit / Bewusstseinstrübung
    {\textless}\,1\Hour*
    \item Übelkeit, Erbrechen, Schwindel
    \item Evtl. flüchtige neurologische Ausfälle
    \item Retrograde  Amnesie  ({\quotedblbase}kann sich an den
    Unfallhergang nicht erinnern{\textquotedblleft} )
      \item Keine  bleibenden Schäden
  \end{compactitem}
  \item \EE{Gehirnprellung -- 2\textdegree}
  \begin{compactitem}
    \item Bewußtlosigkeit / Bewusstseinstrübung {\textgreater}\,1\Hour*
    \item Nachweisbare Gehirnschädigung / bleibende  Ausfälle
    \item Evtl. Gehirnnervenabrisse
    \item  vegetative Störungen
  \end{compactitem}
  \item \EE{Gehirnquetschung -- 3{\textdegree}}
  \begin{compactitem}
    \item Bewußtlosigkeit  {\textgreater}\,24\Hour*
    \item Atmungs-/ Kreislaufstörungen
    \item Hirndrucksymptomatik
    \item Evtl. Einklemmung des Atemzentrums ${\rightarrow}$
    zentraler Atemstillstand
    \item Behinderung der Hirndurchblutung
    \item Hirnblutungen
    %     \begin{itemize}
%         \item Epidural
%         \item Subdural
%         \item Subarachnoidal
%     \end{itemize}
\item Lähmungen
\item Zerebrale Krämpfe
  \end{compactitem}
\end{itemize}
}{}{}{}

\Layout{57}{}{}{}{
\begin{tabularx}{\linewidth}{XXXX}\addlinespace\toprule\addlinespace
\noindent~
&
\noindent\TabHeading{I{\textdegree}}
&
\noindent\TabHeading{II{\textdegree}}
&
\noindent\TabHeading{III{\textdegree}}
\\\addlinespace\midrule\addlinespace
% \TabSub{Grad}
% &
% 1
% &
% 2
% &
% 3
% \\\addlinespace

&
Gerhinerschütterung
&
Prellung
&
Quetschung
\\\addlinespace
\noindent\TabSub{{\ldots} cerebri}
&
Commotio
&
Contusio
&

\\\addlinespace
\noindent\TabSub{Bewusstlosigkeit / Trübung}
&
< 1\Hour*
&
1\Hour* < x < 24\Hour
&
> 24\Hour
\\\addlinespace
\noindent\TabSub{Symptome}
&
Übelkeit, Erbrechen, Schwindel

Evtl. flüchtige neurologische Ausfälle

Retrograde  Amnesie  ({\quotedblbase}kann sich an den
Unfallhergang nicht erinnern{\textquotedblleft} )

Keine  bleibenden Schäden
&
Nachweisbare Gehirnschädigung / bleibende  Ausfälle

Evtl. Gehirnnervenabrisse

vegetative Störungen
&
Atmungs-/ Kreislaufstörungen

Hirndrucksymptomatik

Evtl. Einklemmung des Atemzentrums ${\rightarrow}$
zentraler Atemstillstand

Behinderung der Hirndurchblutung

Hirnblutungen,
Lähmungen, zerebrale Krämpfe

% \\\addlinespace
% Grad
% &
% 
% &
% 
% &

\\\addlinespace\bottomrule
\end{tabularx}
}{SHT-Grade: Übersicht}{}






%%%%%%%%%%%%%%%%%%%%%%%%%%%%%%%%%%%%%%%%%%%%%%%%%%%%%%%%%%%
%%%%%%%%%%%%%%%%%%%%%%%%%%%%%%%%%%%%%%%%%%%%%%%%%%%%%%%%%%%
%%%%%%%%%%%%%%%%%%%%%%%%%%%%%%%%%%%%%%%%%%%%%%%%%%%%%%%%%%%

\clearpage
\section{Wirbelsäulentrauma und Verletzungen des
Rückenmarks}
\label{topic:wirbelsaeulentrauma}

\AuthNotes{KOCH: Überarbeiten, evt etwas kürzen.}
\Text{\TxBeschreibung}
{%
\DefinitionInPara{Ein Wirbelsäulentrauma ist eine
Gewalteinwirkung auf die Wirbelsäule, die zur
Verschiebung, Verrenkung oder zur  Fraktur von Wirbeln mit
oder ohne Rückenmarkschädigung führt \cite{Gorgass7}.}
%
}

\Text{Typische Unfallmechanismen}{
\begin{itemize}
  \item Hochgeschwindigkeitsunfall
  \item Sturz $\geq$ 3fache Patientengröße \A{3fache
  Patientengröße (ITLS) oder 3\Meter*?}
  \item Penetrierende Verletzungen im Bereich der
  Wirbelsäule
  \item Sportverletzungen im Kopf- oder Nackenbereich
  \item Trauma nach Sprung ins Wasser
  \item Bewusstloser Traumapatient
  \item Stauchungsverletzung der \Abk{HWS} (z.\,B. Schlag auf
  Kopf)
  \item Suizidversuch durch Erhängen
\end{itemize}
nach \cite{ITLSde6}, vgl.
\prettyref{MK-m:einschaetzung-indikation-ws-immobilisation} }
{}{}{}{}


%\paragraph{\TxGefahren}
\TextSlide{\TxGefahren}{
\begin{itemize}
  \item Kompression des Rückenmarks (aufgrund von Einblutungen ins  Rückenmark 
  und dadurch bedingtes Ödem)
  \item Durchtrennung des Rückenmarks
  \item Lähmungen (Parese)
  \item Sensibilitätsstörungen (Parästhesie)
  \item Schock (Neurogener / spinaler Schock)
  \item Begleitverletzungen
\end{itemize}
}{
\begin{itemize}
  \item Verletzung des Rückenmarks
  \item Neurologische Ausfälle
  \item Neurogener Schock
  \item Begleitverletzungen
\end{itemize}
}{}{}{}






\TextSlide{Schleudertrauma, Peitschenschlagsyndrom}
{\index{Schleudertrauma}%
\index{Peitschenschlagsyndrom}%
Schleudertraumen kommen bei Auffahrunfällen vor,
insbesondere dann, wenn der Patient  nicht angeschnallt war,
oder die Kopfstütze fehlt. Es kommt zu einer
Zerrung  des Halsmarks, welche sich durch \E{Schmerzen im
\Abk{HWS}-Bereich}, Kopfschmerzen, Schwindel, Tinnitus und
evtl. neurologische Ausfälle (in schweren Fällen)
bemerkbar macht.
}
{
\begin{itemize}
  \item Auffahrunfall
  \item Zerrung des Halsmarks
\end{itemize}
}{}{}{}


\TextSlide{Querschnittssyndrom}{%
Kommt es zu einer Druckerhöhung (z.\,B. aufgrund
einer Blutung) im Rückenmark oder zu einer Durchtrennung des
Rückenmarks können Symptome eines Querschnitts auftreten. Ein
sorgsamer und vorsichtiger Umgang bei der Rettung, der
Untersuchung und dem Transport des Patienten ist unbedingt
notwendig. Es können folgende Symtome auftreten:
\begin{itemize}
  \item Ausfall von Motorik und/oder Sensibilität
  \item Blutdruckabfall aufgrund fehlender Gefäßregulation
  \item Schlaffe oder spastische Lähmung der Extremitäten
  \item Inkontinenz aufgrund von Blasen-/Darmlähmung
  \item Dauererektion \Z{(Priapismus)}
\end{itemize}
}{
\begin{itemize}
  \item Ausfall von Motorik und/oder Sensibilität
  \item Blutdruckabfall
  \item Schlaffe oder spastische Lähmung der Extremitäten
  \item Inkontinenz
  \item Dauererektion
\end{itemize}
}{}{}{}


\TextSlide{Spinaler Schock}
{Aufgrund des Unfalls kommt es zu einer plötzlichen
Durchtrennung des Rückenmarks, wodurch die
Steuerung der Gefäßregulation (und somit der Gefäßtonus)
versagt. Somit sinkt der Blutdruck rasch ab, es kommt zu
einer lebensbedrohlichen Kreislaufstörung. Bei einem hohen
Querschnitt kann auch nervöse die Atemmuskulatur betroffen
sein und es zur Atemlähmung kommen.

\vspace{2\baselineskip}
}{
\begin{itemize}
	\item Querschnittslähmung	
	\item Fehlender Gefäßregulation
\item Atemlähmungen (bei hohem Querschnitt (im Bereich der
\Abk{HWS}))
\end{itemize}
}{}{}{}


% M %%%%%%%%%%%%%%%%%%%%%%%%%%%%%%%%%%%%%%%%%%%%%%%%%%% M %
% Rückenmarksverletzung, Verdacht
\ProcedurePrintDefault{TT09031C}


% M %%%%%%%%%%%%%%%%%%%%%%%%%%%%%%%%%%%%%%%%%%%%%%%%%%% M %
% Rückenmarksverletzung mit Symptomen
\ProcedurePrintDefault{TT09032C}


%%%%%%%%%%%%%%%%%%%%%%%%%%%%%%%%%%%%%%%%%%%%%%%%%%%%%%%%%%%
%%%%%%%%%%%%%%%%%%%%%%%%%%%%%%%%%%%%%%%%%%%%%%%%%%%%%%%%%%%
%%%%%%%%%%%%%%%%%%%%%%%%%%%%%%%%%%%%%%%%%%%%%%%%%%%%%%%%%%%
\clearpage % TEMP temp clearpage
\section{Thoraxtrauma}
\label{topic:thoraxtrauma}

\Text{\TxBeschreibung}
{%
\DefinitionInPara{Ein Thoraxtrauma bezeichnet eine Gewalteinwirkung auf den Brustkorb
mit Verletzung des Brustkorbes bzw. der darin
enthaltenen Organe.}
%
Thoraxverletzungen gelten als besonders gefährlich, da im
Brustkorb (Thorax) lebenswichtige Organe (Herz,
Lunge) und große Gefäße liegen. Fast jeder zweite
Verkehrstote hat ein Thoraxtrauma. \cite{Gorgass7},
\cite{RedelsteinerHRNS1}
}

\Cave{Man darf niemals aufgrund der äußeren Verletzungen auf
die Schwere des Thoraxtraumas schließen, da schwere innere
Verletzungen auftreten können! Oft ist von außen gar keine
Verletzung erkennbar\cite{Strauss0000}.}

%\paragraph{\TxSymptome des Thoraxtraumas}~

\TextSlide{\TxSymptome{} des Thoraxtraumas}
{
Die Hauptgefahr beim Thoraxtrauma ist eine wesentliche
Verletzung des knöchernen Brustkorbes (Sternumfraktur,
Serienrippenfraktur), wodurch ein \EE{instabiler Thorax}
resultiert: Die Atemmechanik ist stark beeinträchtigt, da
sich die Lunge nicht mehr an den Rippen aufspannen
kann. Es kommt folglich zu einer \EE{paradoxen
Atmung}\index{Atmung!paradoxe}, bei der sich der Brustkorb
beim Einatmen
senkt.
Verletzt eine der gebrochenen Rippen die Pleura, oder reisst
diese aufgrund der Gewlteinwirkung ein, ensteht ein
\EE{Pneumothorax} oder ein
Spannungspneumothorax\ZFootnote{Pneumothorax evtl. mit
Hautemphysem (Luftansammlung in der
Haut, typisches Knistern)}.

Der Patient leidet unter \EE{Atemnot} und evtl.
atemabhängigen Schmerzen. Evtl. kommt es zum Bluthusten. Bei
der Inspektion des Brustkorbes können Prellmarken (z.\,B.
von der Lenksäule eines KFZ) auffallen.

Im Brustkorb befinden sich einige lebenswichtige Organe und
Strukturen, daher kann es zu schweren
\EE{Begleitverletzungen} kommen.
Besonders erwähnenswert sind hierbei die Lunge, das Herz,
die Aorta, die Luft- und die Speiseröhre. Auch die Organe
des Oberbauches (Leber, Milz, \ldots) liegen hinter den den
Rippen und können ebenfalls verletzt sein.
}{%
\begin{itemize}
  \item Instabiler Thorax
  \item Evtl. Pneumothorax
  \item Dyspnoe
  \item Atemabhängige Schmerzen
  \item Bluthusten
  \item Prellmarken
  \item Begleitverletzungen!
\end{itemize}
}{}{}{}



\subsection{Pneumothorax}
\label{topic:pneumothorax}%
\index{Pneumothorax}%
\index{Pneumothorax!geschlossener}%
\index{Pneumothorax!offener}%
\index{Pneumothorax!Spontan-}%
\index{Pneumothorax!Spannungs-}%
\index{Spontanpneumothorax}%
\index{Spannungspneumothorax}%

\TextSlide{\TxBeschreibung{} und Arten}{%
\DefinitionInPara{Ein Pneumothorax ist eine Luftansammlung im Pleuraspalt.}
%
Zwischen Lungenoberfläche und innerer Thoraxwand herrscht
beim Gesunden ein Unterdruck gegenüber dem äußeren Luftdruck.
Die Lunge ist dadurch flexibel im Brustkorb aufgespannt. Beim
Pneumothorax gelangt durch eine Verletzung Luft in diesen
Pleuraspalt, welche den Unterdruck aufhebt und die Lunge
zieht sich zusammen. 
Vgl. Anatomie: \FullPrettyRef{topic:atemmechanik}, und
\FullPrettyRef{topic:pleura}.

Man unterscheidet zwischen \T{offenem
Pneumothorax} (Luft strömt durch eine Verletzung der
Brustwand von außen in den Pleuraspalt) und
\T{geschlossenem Pneumothorax} (Luft strömt durch eine
Verletzung der Lungenoberfläche von innen in den
Pleuraspalt). 

Entsteht durch die Form der Wunde ein Ventilmechanismus wird
beim Einatmen Luft in den Pleuraspalt gesogen, während sich
beim Ausatmen die Wunde verschließt und die Luft nicht mehr
entweichen kann. Dadurch füllt sich der Pleuraspalt mehr und
mehr mit Luft, bis schließlich im betroffenen Pleuraspalt
ein Überdruck entsteht. Im Extremfall drückt der Überdruck
den Mittelfellraum mit dem Herz auf die Gegenseite, wodurch
auch noch die gesunde Lungenhälfte und evtl. die Hohlvene
komprimiert wird. Im schlimmsten Fall kann das Herz an
seiner Bewegung gehindert werden. Tritt bei einem
Pneumothorax ein derartiger Ventilmechanismus auf, spricht
man vom \T{Spannungspneumothorax}.

Ein Pneumothorax kann auch ohne vorangegangene Verletzung
auftreten. Diese Form des Pneumothorax nennt man
\T{Spontanpneumothorax} und tritt am häufigsten bei jungen,
hochgewachsenen und schlanken Menschen auf. Symptome des
Spontanpneumothorax sind Schmerzen, Atemnot und Todesangst.
}{
\begin{itemize}
  \item Offener Pneumothorax
  \item Geschlossener Pneumothorax
  \item Spannungspneumothorax
  \item Spontanpneumothorax
\end{itemize}
}{}{}{}




\begin{PictureSeriesWide}{}{Bilderserie \SlideHeading{Pneumothorax}}
\PictureSeriesRowWide{3}
{./images/pneumothorax_xray_halb.jpg}
{Pneumothorax. 
Beachte die Randlinie zwischen Lunge
(unten, innen) und dem luftgefüllten
Pleuraspalt (oben, außen) \SourcePicture{Hauer}}
{./images/pneumothorax_xray_endstadium.jpg}
{Kompletter
Pneumothorax der rechten Seite. Die Lunge ist kollabiert.  
\SourcePicture{Hauer}}
{./images/pneumothorax_stichverletzung_diskret.jpg}
{Diese
unauffällige Stichverletzung hat einen Pneumothorax
verursacht.
\Ca{Dieser Patient ist lebensgefährlich verletzt!} 
\SourcePicture{Hauer}}
{}
{}
\end{PictureSeriesWide}


% M %%%%%%%%%%%%%%%%%%%%%%%%%%%%%%%%%%%%%%%%%%%%%%%%%%% M %
% Pneumothorax
\ProcedurePrintDefault{TJ93090C}


%\AuthNotes{KOCH: 

%Es fehlt:
%Serienrippenfraktur erklären
%Paradoxe Atmung
%Gefahr von Blutverlust
%Therapie
%Evtl. auch kurz etwas über STernumfraktur schreiben?}

\subsection{Rippenfraktur, Serienrippenfraktur und instabiler Thorax}
\label{topic:serienrippenfraktur}
\index{Serienrippenfraktur}

\TextSlide{\TxBeschreibung}
{%
Von einer \EE{Serienrippenfraktur} spricht man,
wenn mindestens drei benachbarte Rippen gebrochen sind. Sind
die Rippen an verschiedenen Stellen gebrochen oder gar
zertrümmert, werden diese Teile bei der Einatmung nach innen
gezogen und bei der Ausatmung nach außen gedrückt. Man
bezeichnet dies (Brustkorb \E{senkt} sich bei der Einatmung)
als \II[Atmung!paradoxe!bei Serienrippenfraktur]{paradoxe
Atmung} bei \EE{instabilem Thorax}. Die Folge ist eine
hochgradige Atemnot, außerdem treten (z.\,T. heftige)
Schmerzen auf, die sich bei der Einatmung verstärken. Rippenbruchstücke können
darunter liegende Organe verletzen. Es kann
zu Lungeneinrissen, sowie Leber- und Milzverletzungen mit
starken inneren Blutungen kommen. Der Blutverlust kann \VonBis{2}{3} Liter
betragen \cite{LPN3}. In diesem Fall herrscht höchste
\E{Schockgefahr}.
}{
\begin{itemize}
  \item Schmerzen beim Einatmen
  \item Evtl. Atemnot, Zyanose, Tachypnoe
  \item Evtl. paradoxe Atmung
  \item Evtl. Schocksymptome
\end{itemize}
}{}{}{}

\subsection{Sternumfraktur}\label{topic:sternumfraktur}\index{Sternumfraktur}
Eine Fraktur des Brustbeins (Sternum) tritt vor allem bei
nicht angeschnallten Fahrzeuglenkern auf, wenn der Brustkorb
gegen das Lenkrad prallt. Meistens treten starke Schmerzen
auf, gelegentlich lassen sich eine Delle oder eine Stufe
tasten. \cite{LPN3}

\subsection{Verletzung des Herzens und der großen Gefäße}

\TextSlide{\TxBeschreibung}{
Verletzungen des Herzens und der großen Gefäße treten am
wahrscheinlichsten bei Stürzen aus großer Höhe und
Verkehrsunfällen mit Frontalaufprall und großer
Geschwindigkeit auf, seltener bei Schuss- oder
Stichverletzungen. Folgende Verletzungen sind typisch:

\begin{itemize}
  \item \EE{Aortenruptur}: Durch einen Einriss der Aorta
  kommt es zu schweren inneren Blutungen. Siehe auch
  Abschnitt \ref{topic:aortenruptur}.
  \item \EE{Herzkontusion}: Hier kommt es durch Quetschungen
  der Herzmuskulatur zu Sickerblutungen. Bei
  Sternumfrakturen mit Herzrhythmusstörungen sollte man an
  eine Herzkontusion denken.
  \item \EE{Herzbeuteltamponade}: Es kommt zu einer
  Einblutung in den Herzbeutel,  wodurch das Herz
  komprimiert wird und sich nicht mehr mit Blut füllen 
  kann. Leitsymptome sind gestaute Halsvenen und schwerer
  Schockzustand (kardiogener Schock).
\end{itemize}
}{
\begin{itemize}
  \item Aortenruptur
  \item Herzkontusion
  \item Herzbeuteltamponade
\end{itemize}
}{}{}{}

\clearpage % TEMP temp clearpage
\section{Bauch- und Beckentrauma}
\index{Bauchtrauma}
\label{topic:bauchtrauma}

\TextSlide{\TxBeschreibung}{
Ein Bauchtrauma oder Abdominaltrauma ist eine Verletzung des
Abdomens. In einigen Büchern ist das Bauchtrauma als
Untergruppe des akuten Abdomens klassifiziert.
Häufigste Ursachen sind Verkehrsunfälle, Stürze aus großer
Höhe sowie Stoß- und Schlagverletzungen.
\cite{RedelsteinerHRNS1} Grundsätzlich ist zwischen offenem
(perforierendem, penetrierendem) und geschlossenem
(stumpfem) Bauchtrauma zu unterscheiden. Alle Bauchtrauma
können mit massiven inneren
Blutungen verbunden sein. }{
\begin{itemize}
  \item Offenes (perforierendes) Bauchtrauma
  \item Geschlossenes (stumpfes) Bauchtrauma
\end{itemize}
}{}{}{}

\Layout{57}{}
{}
{}
{\begin{tabularx}{\linewidth}{XX}
\TabHeading{Offenes (perforierendes) Bauchtrauma}
&
\TabHeading{Geschlossenes (stumpfes) Bauchtrauma}
\\\midrule
\begin{itemize}
  \item  Eröffnung der Bauchhöhle, evtl. Austritt von Darmschlingen
  \item Stich-, Pfählungsverletzung
  \item Schnittverletzung, {\quotedblbase}Platzbauch{\textquotedblleft}
\end{itemize}
&
\begin{itemize}
  \item Schlag, Tritt
  \item Prellung durch Lenkstange bei PKW-Unfall
\end{itemize}
\\\bottomrule
\end{tabularx}}
{Offenes und geschlossenes Bauchtrauma}
{}


\subsection{Offenes Bauchtrauma}
\label{topic:bauchtrauma-offen}

\Layout{20}{\TxBeschreibung}{%
Offene Bauchtraumata sind seltener als geschlossene
Bauchtraumata. Beim offenen Bauchtrauma kann es zum Austritt
von Darmschlingen kommen. Die Größe der Eintrittswunde sagt
nichts über den Schweregrad der Verletzung aus, da unter der
Eintrittswunde unterschiedlich schwere Organschäden mit
inneren Blutungen versteckt sein können. \Z{In 25\,\% der
Fälle treten Begleitverletzungen im Brustkorb auf
\cite{Strauss0000}.}}{
\begin{itemize}
  \item Austritt von Darmschlingen
  \item unterschiedlich schwere Organschäden
  \item innere Blutungen
\end{itemize}
}
{./images/trauma/bauchtrauma-offenes-01.jpg}
{Offenes Bauchtrauma mit Austritt von Darmschlingen.}
{Hauer}


% M %%%%%%%%%%%%%%%%%%%%%%%%%%%%%%%%%%%%%%%%%%%%%%%%%%% M %
% Bauchtrauma, offen
\ProcedurePrintDefault{TS31011C}



\subsection{Geschlossenes Bauchtrauma}
\label{topic:bauchtrauma-geschlossen}

Geschlossene Bauchtraumata sind aufgrund der inneren
Blutungen besonders gefährlich. So kann es z.\,B. bei einer
Milzruptur zu einem Blutverlust von 4~Liter kommen!
\cite{Gorgass7}

\Fazit{Immer nach \EE{Prellmarken} suchen. \EE{Patienten
entkleiden!}}

\TextSlide{\TxSymptome}
{%
\begin{itemize}
  \item Starke  Bauchschmerzen
  \item Brettharte  Bauchdecke  (Abwehrspannung, D\'efense )
  \item Leitsymptom: Schockzeichen
  \item Beim Bewußtlosen schwierig zu  diagnostizieren
  \begin{itemize}
    \item Gründlich nach \index{Prellmarken}\EE{Prellmarken}
    suchen!
  \end{itemize}
\end{itemize}
}{
\begin{itemize}
  \item Brettharte  Bauchdecke
  \item Prellmarken!
\end{itemize}
}{}{}{}


\TextSlide{\TxGefahren}{%
\begin{itemize}
  \item Gefäßverletzungen
  \begin{itemize}
    \item Aorta, Vena cava, \ldots{}
    \item Massive innere Blutungen ${\rightarrow}$ Volumenmangelschock
  \end{itemize}
  \item Zerreißung  von Organen (\EE{Milz},
  \EE{Leber})
  \begin{itemize}
    \item Massive \EE{innere Blutungen} (bis über 4 Liter) ${\rightarrow}$
    Volumenmangelschock
  \end{itemize}
  \item Zerreißung  von Hohlorganen (Magen, Darm,Gallengänge )
  \begin{itemize}
    \item Austritt von Speisebrei, Magensaft, Stuhl, Gallensekret
    \begin{itemize}
      \item Bauchfellentzündung (Peritonitis) nach Std. bis  Tagen
      \item Sepsis (hohe Letalität), septischer Schock
    \end{itemize}
  \end{itemize}
\end{itemize}
}{
\begin{itemize}
  \item Gefäßverletzungen
  \item Zerreißung  von Organen
  \item Zerreißung  von Hohlorganen
\end{itemize}
}{}{}{}








\subsubsection{Milzruptur}
\index{Milzruptur}
\label{topic:milzruptur}

\TextSlide{\TxBeschreibung{} und \TxGefahren{}}{%
Die Milz ist ein \enquote{teuflisches} Organ. Sie besitzt
rund um das gut durchblutete Gewebe eine Bindegewebskapsel.
Das Schlimme daran: Das Gewebe kann reißen, die Kapsel
bleibt jedoch intakt. Es kommt zuerst nur zu einer
geringfügigen Einblutung, dem Patienten geht es zunächst
gut. Nach einiger Zeit reißt dann die Kapsel und es kommt zu
einer ungehemmten Blutung in den Bauchraum. \EE{Der Patient
verfällt von einem Moment auf den anderen in einen
Volumenmangelschock.}

\begin{itemize}
  \item \EE{Primäre (einzeitige) Milzruptur}: Die Milz und
  die Kapsel reißen gleichzeitig, es kommt zu einem unmittelbaren
  Schockzustand.
  \item \EE{Sekundäre (zweizeitige) Milzruptur}:
  \label{topic:milzruptur-zweizeitig} Die Kapsel reißt erst
  einige Zeit (Stunden bis Tage) nach dem Riss des Milzgewebes
  \E{(Latenzzeit)}. Es kommt zu einem \EE{zeitlich versetzten} und oft
  \EE{unerwarteten Schockgeschehen}. Mitunter reißt die
  Kapsel erst Stunden nach dem eigentlichen Unfall.
  
  \E{Daher
  ist es wichtig jeden Patienten, bei dem auch nur der geringste
  begründete Verdacht auf eine Milzverletzung besteht, einer
  weiterführenden Untersuchung im Spital
  zuzuführen!}\ZFootnote{Mittels Ultraschall oder
  Computertomographie (CT) kann eine Milzblutung
  nachgewiesen werden.}
\end{itemize}
}{
\begin{itemize}
  \item Schmerzen im linken Oberbauch
  \item Manchmal Ausstrahlung der Schmerzen in die linke Schulter (Kehr-Zeichen)
  \item Volumenmangelschock
  \item Möglicherweise zeitversetzter Kapselriss
  (sekundäre Milzruptur, unerwartetes Schockgeschehen)
  \item Primäre Milzruptur
  \item Sekundäre Milzruptur:
  \begin{itemize}
    \item Erst nur Riss des Milzgewebes, Kapsel intakt
    \item Symptomfreies Intervall (Latenzzeit), Sickerblutung
    \item Später reißt auch Kapsel  ${\rightarrow}$
    plötzlich starke Blutung
  \end{itemize}
\end{itemize}
}{}{}{}



\Cave{Cave! {\quotedblbase}\index{Zweizeitige
Milzruptur}Zweizeitige Milzruptur{\textquotedblleft}}






\subsubsection{Weitere innere Verletzungen}
\index{Leberruptur}\label{topic:leberruptur}
\index{Aortenruptur}\label{topic:aortenruptur}
\index{Nierenruptur}\label{topic:nierenruptur}

\TextSlide{\TxBeschreibung}{
Der Bauchraum beherbergt viele, teilweise sehr gut
durchblutete Organe die bei einer Verletzung zu erheblichen
Blutverlusten führen können. Im Folgenden werden die
Leberruptur, die Aortenruptur sowie die Nierenruptur näher
beschrieben. Präklinisch äußern sich diese Verletzungen im
durch den Blutverlust hervorgerufenen Schockzustand. Weitere
Hinweiszeichen sind Prellmarken oder Schmerzen in der
entsprechenden Gegend.

Außerhalb eines Spitals kann man derartige Verletzungen kaum
ursächlich behandeln, im Vordergrund steht die Erkennung,
Schockbehandlung und der zügige Transport in ein
Traumazentrum.
}{
\begin{itemize}
  \item Leberruptur
  \item Aortenruptur
  \item Nierenruptur
\end{itemize}
}{}{}{}


\TextSlide{Leberruptur}
{%
Die Leber ist auch gut durchblutet und kann im Falle einer
Verletzung große Probleme machen. Ursachen sind meist
stumpfe Verletzungen, z.\,B. ein Lenkradaufprall oder eine
Sicherheitsgurtkompression. Der Patient klagt meist über
einen starken Drucksschmerz im rechten Oberbauch. Durch den
Blutverlust durch die Verletzung dieses stark durchbluteten
Organs kann es zu einem Volumenmangelschock kommen. Wie bei
der Milz kann es auch hier zu einer zweizeitigen Leberruptur
kommen.
}{
\begin{itemize}
  \item Druckschmerz im rechten Oberbauch
  \item Volumenmangelschock
  \item Späterer Kapselriss möglich (zweizeitige Ruptur wie
  bei Milz, siehe oben)
\end{itemize}
}{}{}{}

\TextSlide{Aortenruptur}{
Durch
einwirkende starke Kräfte kann es zu einem Riss in der Aorta
kommen. Besonders bei Stürzen aus großen Höhen kann dies
vorkommen. Da es sich bei der Aorta um das größte Gefäß des
Körpers handelt ist die Prognose äußerst schlecht.
}{
\begin{itemize}
  \item Riss in der Aorta
  \item schlechte Prognose
\end{itemize}
}{}{}{}

\Text{Nierenruptur}{
Die Niere als naturgemäß gut durchblutetes Organ kann im
Falle einer Verletzung auch massive Probleme machen.
}{

}{}{}{}

\subsection{Beckentrauma}
\index{Beckentrauma}
\label{topic:beckentrauma}

\TextSlide{\TxBeschreibung{} und \TxSymptome}{
Ein Beckentrauma ist eine Verletzung des Beckens. In manchen
Lehrbüchern zählt es zu den Bauchverletzungen.

Besonders gefährlich sind Frakturen des Beckengürtels
(\prettyref{topic:beckenguertel}). Stehen und Gehen ist dann
nicht mehr möglich, der Patient verspürt starke Schmerzen
bei jeder passiven Bewegung \E{beider} Beine. Da der
Beckenknochen gut durchblutet ist, können binnen weniger
Minuten bis zu fünf Liter Blut (also fast die gesamte
Blutmenge) verloren gehen. Es besteht höchste
\EE{Schockgefahr!} Eine \EE{Kompression des Beckens} ist \EE{überlebenswichtig}
um den möglichen großen Blutverlust einzudämmen.

\begin{itemize}
  \item Entstehung durch erhebliche Gewalteinwirkung 
  (Verschüttung, Überrolltrauma)
  \item Massive  Blutungen (bis 5\Liter* in wenigen  Minuten),
  Gefahr des Ausblutens (evtl. ohne sichtbare  Blutungsquelle!)
  \item Mitverletzung von Organen im kleinen  Becken:
  \begin{itemize}
    \item Dickdarm ${\rightarrow}$ Blutung aus Anus
    \item Harnsystem ${\rightarrow}$ Blut im Harn, Harn in der
    Bauchhöhle
    %\item (Ausnahme: Spontanruptur der gefüllten
    %Harnblase bei stumpfem Bauchtrauma)
    \item Hoden ${\rightarrow}$ Hämatome am Skrotum und am  Damm
  \end{itemize}
\end{itemize}
}{
\begin{itemize}
  \item Verletzung des Beckens
  \item Massive innere Blutungen $\rightarrow$ Schockzeichen
  \item Kompression des Beckens ist überlebenswichtig!
  \item Verlust von bis zu 5 Litern Blut (= gesamtes
  Blutvolumen des Menschen)
  möglich!
\end{itemize}
}{}{}{}


% M %%%%%%%%%%%%%%%%%%%%%%%%%%%%%%%%%%%%%%%%%%%%%%%%%%% M %
% Beckentrauma, instabil
\ProcedurePrintDefault{TS32831C}


%%%%%%%%%%%%%%%%%%%%%%%%%%%%%%%%%%%%%%%%%%%%%%%%%%%%%%%%%%%
%%%%%%%%%%%%%%%%%%%%%%%%%%%%%%%%%%%%%%%%%%%%%%%%%%%%%%%%%%%
%%%%%%%%%%%%%%%%%%%%%%%%%%%%%%%%%%%%%%%%%%%%%%%%%%%%%%%%%%%
\clearpage % TEMP temp clearpage
\section{Polytrauma}
\label{topic:polytrauma}

\Text{\TxBeschreibung}
{%
\DefinitionInPara{Der Begriff \T{Polytrauma} bezeichnet
gleichzeitige \EE{Verletzungen mehrerer Körperregionen} oder
Organsysteme, von denen \EE{wenigstens eine} allein
\EE{oder die Kombination mehrerer} \EE{lebensgefährlich}
ist \cite{BoehmerTaschRett6,Gorgass7,RedelsteinerHRNS1,LPN3}. 
}
%
Die Arbeitsdiagnose Polytrauma kann auch ausschließlich
aufgrund eines Verdachts, begründet auf den Unfallmechanismus und eventuell der Vitalparameter,
gestellt werden.

Beispiele:
\begin{itemize}
  \item Bauchtrauma + Oberschenkel; \Abk{SHT} + WS + Thoraxtrauma etc.
  \item Gebrochener Finger + gebrochene Zehe ${\rightarrow}$ kein
  Polytrauma
\end{itemize}

Häufigkeit verletzter Körperregionen beim Polytrauma
\cite{BoehmerTaschRett6,Gorgass7}:
\begin{itemize}
  \item \VonBis{60}{67}\PC* Schädel-Hirn
  \item \VonBis{45}{75}\PC* Extremitäten
  \item \VonBis{25}{30}\PC* Thorax
  \item \VonBis{10}{37}\PC* Abdomen und Becken
  \item \VonBis{5}{15}\PC* Wirbelsäule
\end{itemize}

Die Arbeitsdiagnose Polytrauma kann auch ausschließlich
aufgrund des Unfallmechanismus und der Vitalparameter
gestellt werden.

\Fazit{Der Patient gilt als 
polytraumatisiert bis zum Beweis des Gegenteils! Jedes
Polytrauma stellt eine \EE{vitale Bedrohung} dar. }

Die Sterblichkeit bei polytraumatisierten Patienten liegt
zwischen \VonBis{20}{50}\PC*, ist also hoch. Ein gezieltes,
strukturiertes und zügiges Vorgehen ist daher unabdingbar!
Im Idealfall sollten zwischen dem Eintreffen der
Rettungskräfte und der definitiven Versorgung im Spital
nicht mehr als 60 Minuten vergehen \cite{BoehmerTaschRett6}.
Die Behandlung muss sich daher in erster Linie an der
Stabilisierung der Vitalfunktionen und nicht an den
Einzelverletzungen orientieren!

\Fazit{Ein polytraumatisierter Patient ist ein
zeitkritischer Patient!}


% M %%%%%%%%%%%%%%%%%%%%%%%%%%%%%%%%%%%%%%%%%%%%%%%%%%% M %
% Polytrauma, allgemein
\ProcedurePrintDefault{TT07XX1C}



%%%%%%%%%%%%%%%%%%%%%%%%%%%%%%%%%%%%%%%%%%%%%%%%%%%%%%%%%%%
%%%%%%%%%%%%%%%%%%%%%%%%%%%%%%%%%%%%%%%%%%%%%%%%%%%%%%%%%%%
%%%%%%%%%%%%%%%%%%%%%%%%%%%%%%%%%%%%%%%%%%%%%%%%%%%%%%%%%%%
\clearpage % TEMP temp clearpage
\section{Extremitätentrauma und Sportverletzungen}
\label{topic:extrimitaetentrauma}
%Gelenke betreffend
Verletzungen der Extremitäten können in verschiedenen Formen auftreten:
\begin{itemize}
  \item Verstauchung (Distorsion, \FullPrettyRef{topic:verstauchung-tra})
  \item Verrenkung (Luxation, \FullPrettyRef{topic:verrenkung})
  \item Knochenbruch (Fraktur, \FullPrettyRef{topic:frakturen})
  \item Amputationen
\end{itemize}



\subsection{Amputationen}

\Text{\TxBeschreibung}
{%
\DefinitionInPara{Eine \T{Amputation} 
(\Lang{Lat.} \I{Amputatio} (\I{Amp.})) 
bezeichnet einen kompletten oder
inkompletten  Abriss bzw. Abtrennung (von Teilen) einer
Extremität oder eines anderen Körperteils, sodass deren
Durchblutung ganz oder teilweise aufgehoben ist. Den
abgetrennten Körperteil nennt man \T{Amputat}.
\cite{BoehmerTaschRett6}}
%
Amputationen erfolgen meist bei Motorradunfällen oder
Schnittverletzungen (z.\,B. mit einer Kreissäge). Je glatter
die Wundränder sind desto größer ist die Chance dass das
Amputat wieder erfolgreich angenäht werden kann. Nach dem
Amputat ist in jedem Fall zu suchen. Der Amputatstumpf wird
steril mit einem Druckverband versorgt (kein Abbinden!). An
große Blutungen und Schockgefahr denken!
}



% M %%%%%%%%%%%%%%%%%%%%%%%%%%%%%%%%%%%%%%%%%%%%%%%%%%% M %
% Versorgung eines Amputats
\ProcedurePrintDefault{TT14071C}


\subsection{Weitere Verletzungsbilder}

\TextSlide{\TxBeschreibung}{%
Neben den bereits genannten Verletzungen, treten (im Sport
oder in der Freizeit) oft auch folgende, meist harmlose,
Verletzungen auf, welche hier der Vollständigkeit halber
erwähnt sind:

\begin{itemize}
  % 	\item Muskelkrampf
     \item Muskelzerrung
     \item Muskelfaserriss
     \item Achillessehnenriss
     \item Meniskusschaden
\end{itemize}
}{
\begin{itemize}
  %  \item Muskelkrampf
     \item Muskelzerrung
     \item Muskelfaserriss
     \item Achillessehnenriss
     \item Meniskusschaden
\end{itemize}
}{}{}{}

%%%%%%%%%%%%%%%%%%%%%%%%%%%%%%%%%%%%%%%%%%%%%%%%%%%%%%%%%%%
%%%%%%%%%%%%%%%%%%%%%%%%%%%%%%%%%%%%%%%%%%%%%%%%%%%%%%%%%%%
%%%%%%%%%%%%%%%%%%%%%%%%%%%%%%%%%%%%%%%%%%%%%%%%%%%%%%%%%%%
%%%%%%%%%%%%%%%%%%%%%%%%%%%%%%%%%%%%%%%%%%%%%%%%%%%%%%%%%%%
\clearpage % TEMP temp clearpage
\section{Verbrennung und Verbrühung}
\label{topic:verbrennung}%

\TextSlide{\TxBeschreibung}{%
\DefinitionInPara{Eine Verbrennung 
\Z{(\Lang{Lat.} \I{Combustio}, \Lang{Abkz.} \I{Comb.})} 
ist eine durch Hitzeeinwirkung
entstandene Schädigung des Gewebes, besonders häufig der
Haut. Sie ist eine thermische
Wunde.} 
\DefinitionInPara{Eine \T{Verbrühung} 
\Z{(\Lang{Lat.} \I{Ambustio} (\I{Amb.}), \Lang{engl.} scald)} 
ist eine thermische Schädigung durch heiße Flüssigkeit oder Dampf
\cite{Pschy259}.}%
\noindent\Commentary[Verbrennung / ASBÖ-Lehrmeinung]
{Dieser Abschnitt ist kompatibel mit der ASBÖ-Lehrmeinung,
publiziert unter 
\Speech{\protect\bibentry{ASBOEAkademieUpdate2010Verbrennung}}
\cite{ASBOEAkademieUpdate2010Verbrennung}.}
%
Verbrennungen reichen von kleinen Brandblasen bis schweren
Schädigungen der Haut und der darunter liegenden Gewebe
(z.\,B. großflächige Verkohlung). Um die Schwere der
Verletzung rasch einschätzen zu können müssen die Ausdehnung
und der Schweregrad der Verbrennung beurteilt werden. Dazu
stehen die \E{Neunerregel} sowie die Einteilung nach
(Schwere-)Graden zur Verfügung. Schwere Verbrennungen können
-- wie alle Wunden -- Auswirkungen auf den gesamten
Organismus haben.%
\Commentary[Verbrennung / Unterscheidung leichte
und schwere Verbrennung]
{Um die
Unterscheidung zwischen leichter und schwerer Vebrennung
für den Leser einfacher zu gestalten, wurden diese
Unterscheidung hier eingeführt und die Begriffe im Konsens
und nach Maßgabe der einschlägigen Literatur entsprechend
definiert.}

Abhängig von den Umständen muss auch an ein
\E{Inhalationstrauma} (\prettyref{topic:inhalationstrauma})
oder andere Begleitverletzungen gedacht
werden.
}{
\begin{itemize}
  \item Verbrennung
  \item Verbrühung
  \item Verschiedene Schweregrade
  \item Neunerregel
  \item Cave: Inhalationstrauma
\end{itemize}
}{}{}{}


\TextSlide{\TxGefahren}{%
In den verbrannten Arealen \EE{verliert die Haut ihre
Schutzfunktion}; es besteht die Gefahr von \E{Verdunstung},
\E{Unterkühlung} und \E{Infektionen}.
\index{Verbrennungskrankheit}
\label{topic:verbrennungskrankheit}
Ab ca. \VonBis{20}{30}\PC* verbrannter Körperoberfläche kann
sich eine \T{Verbrennungskrankheit} entwickeln. Dabei kommt
es zu einer körperweiten \EE{Entzündungsreaktion}. Besonders
dramatisch ist die dadurch entstehende erhöhte
Durchlässigkeit der Blutgefäße: Dadurch tritt
Plasmaflüssigkeit in das Gewebe über und es bilden sich
Ödeme. 
Aufgrund des \EE{Flüssigkeitsverlustes} durch die
Verbrennungskrankheit und Verdunstung besteht bei
großflächigen Verbrennungen \EE{Schockgefahr}
(hypovolämischer Schock)! Der Schock entwickelt sich
allerdings erst allmählich. 

Bei \EE{ringförmige Verbrennungenn}, die um den Brustkorb
herumlaufen, kann sich das Gewebe zusammenziehen und die
Dehnbarkeit des Brustkorbes abnehmen; dadurch wird die
Atmung stark behindert.
\EE{Unterkühlung} bedeutet erhöhte
\E{Sterblichkeit}.\ZFootnote{\cite{Loennecker2001} spricht
bei Verbrennungspatienten von einer über 40\PC*ig erhöhten
Chance zu versterben wenn die Körpertemperatur um 1 Grad
abnimmt. Eine Unterkühlung tritt vor allem bei
narkotisierten Patienten leicht ein.}

}{
\begin{itemize}
  \item Haut verliert Schutzfunktion
  \item Verbrennungskrankheit, Schock
  \item Ringförmige Brustkorbverbrennung: Atembehinderung
  \item Unterkühlung
\end{itemize}
}{}{}{}


\TextSlide{Ausdehnung: Handregel und Neunerregel}
{%
Das Verbrennungsausmaß, also die Größe des geschädigten
Bereichs, wird in Prozent der verbrannten Körperoberfläche
angegeben. Diese kann entweder mit der \E{Handregel} oder
mittels der
 \E{Neunerregel} geschätzt werden. Die \T{Handregel} besagt,
 dass die Handfläche (ohne Finger!)
des Patienten zirka einem Prozent seiner Körperoberfläche
entspricht.

Bei großflächigen Verbrennungen wird die \T{Neunerregel}
angewendet: Hier wird der Körper in 9\PC*-Teile eingeteilt
um das Ausmaß der Vebrennung abschätzen zu können.


\TODO{GABS}{2010-04-18}{Bild: Neunerregel}{}
}
{\begin{itemize}
    \item Handregel
    \item Neunerregel
\end{itemize}}
{}
% {./icons/under-construction.png}
{}{}

\begin{table}[H]
\caption{Neunerregel}\label{tab:Neunerregel}
\begin{tabulary}{\linewidth}[H]{LCC}
\addlinespace\toprule\addlinespace
\noindent & \noindent\TabHeading{\PC} &\noindent 
\\\addlinespace\midrule\addlinespace
Kopf:
 & 9  &
\\
Arm, je:  &  9  &
\\
Rumpf vorne: & 18  &
\\
Rumpf hinten: & 18 &
\\
Bein, je: & 18 &
\\
Genitalien:  & 1 &
\\\addlinespace\bottomrule
\end{tabulary}
\end{table}

\TextSlide{Grade}{%
\index{Verbrennungspanzer}%
Die Schweregrade einer Verbrennung werden wie folgt eingeteilt:

\begin{center}
\begin{tabulary}{\linewidth}[H]{CL}
\addlinespace\toprule\addlinespace
\noindent\TabHeading{Grad}
&
\noindent\TabHeading{Symptome}
\\\addlinespace\midrule\addlinespace
1\textdegree
 &
Rötung, Schmerz, Schwellung
\\\addlinespace
2\textdegree
 &
Blasenbildung, Schmerz
\\\addlinespace
3\textdegree
 &
\T{Verschorfung} (Vebrennungspanzer), Schmerzen
abhängig von der Tiefe der Schädigung
(Zerstörung von Nerven), Schrumpf\-ung des Gewebes durch
Hitze\-ein\-wirkung, reicht bis zur Verkohlung 
\\\addlinespace\bottomrule
\end{tabulary}
\end{center}
}{
\begin{itemize}
  \item 1\textdegree : Rötung, Schwellung, Schmerzen
  \item 2\textdegree : Blasenbildung, Schmerzen
  \item 3\textdegree : Schorfbildung, Verbrennungspanzer
\end{itemize}
}{}{}{}

\begin{PictureSeriesWide}{}{Bilderserie: \EE{Verbrennungen}}
\PictureSeriesRowWide{3}
{./images/trauma/combustio-sa-01.jpg}
{Großflächige
Verbrennung \SourcePicture{Hauer}}
{./images/trauma/combustio-sa-02-hand.jpg}
{Ablederung
\SourcePicture{Hauer}}
{./images/trauma/combustio-sa-03-escharotomie.jpg}
{Escharatomie.  Durch die
Ringförmige Verbrennung am Thorax zieht sich die Haut
zusammen  und die Atmung erschwert (Verbrennungspanzer).
Durch die
{Entlastungsschnitte} kann der Patient wieder atmen.
\SourcePicture{Hauer}}
{}
{}
\end{PictureSeriesWide}





\TextSlide{Einschätzung}{%
Man unterscheidet zwischen einer \E{leichten} und einer
\E{schweren} Verbrennung, maßgebliche Faktoren sind die
\E{Ausbreitung} und der \E{Verbrennungsgrad}. Ab einer
betroffenen Körperoberfläche von >\,20\PC* (Erwachsene;
Kinder >\,10\PC*, Säuglinge >\,5\PC*) gilt eine Verbrennung
als \E{großflächig}. Aufgrund des \E{Flüssigkeitsverlustes}
durch die Verbrennungskrankheit und Verdunstung besteht
bereits ab ab \EE{10\PC*} (oder \E{5\PC* bei Kindern})
\EE{$\geq$\,2\,{\textdegree}} verbrannter Körperoberfläche
\EE{Schockgefahr} (hypovolämischer Schock), diese
Verbrennungen gelten daher als \EE{schwer}! Der Schock
entwickelt sich allerdings erst allmählich.

\Fazit{Ab \EE{10\PC* $\geq$\,2\,{\textdegree}} verbrannter
Körperoberfläche gilt eine Verbrennung als \enquote{schwer}, für
\EE{Kinder} gilt eine Grenze von
\EE{5\PC*.}}

}{Ab \EE{10\PC* $\geq$\,2\,{\textdegree}} verbrannter
Körperoberfläche gilt eine Verbrennung als \enquote{schwer}; für
\EE{Kinder} gilt eine Grenze von
\EE{5\PC*}.
}{}{}{}

% TODO INHALT: Verbrennung:, leichte, Maßnahmen
% Schockbekämpfung?
% Notarzt? wann leichte Verbennung?


% M %%%%%%%%%%%%%%%%%%%%%%%%%%%%%%%%%%%%%%%%%%%%%%%%%%% M %
% Verbrennung, leicht
\ProcedurePrintDefault{TT30001C}


% M %%%%%%%%%%%%%%%%%%%%%%%%%%%%%%%%%%%%%%%%%%%%%%%%%%% M %
% Verbrennung, schwer
\ProcedurePrintDefault{TT30002C}


\Text{Was man \Ca{nicht} tun darf \ldots}
{%
\begin{itemize}
    \item[\Irritant] NEIN: Salbe, Puder, etc. auf frische 
    Brandwunde \cite {emergency9,PHTLS6}
    \item[\Irritant] NEIN: Blasen öffnen \cite{emergency9}
    \item[\Irritant] NEIN: festklebende Kleidung gewaltsam entfernen
%     \item[\Irritant] NEIN: aluminiumbeschichteten Verband  verwenden (Gefahr der
%     Hitzereflexion bei  nicht ausreichend gekühlten 
%     Verbrennungen!)\ACh{GABS}{2010-08-06}{Metallisierte Verbände sind
%     wünschenswert, werden zB in Helfen2008 empfohlen.}
\end{itemize}
}{

}{}{}{}


% \Text{Literatur}{%
\cite{Kamolz2009,Verbrennungsmed2009Praehosp,Beneker2004,Loennecker2001,TertialeNotfallIntensivmed1}
% }{}{}{}{}



% \clearpage % TEMP temp clearpage
\subsection{Inhalationstrauma}
\label{topic:inhalationstrauma}


\TextSlide{{\TxBeschreibung} und {\TxSymptome}}
{%
\DefinitionInPara{Von einem \T{Inhalationstrauma} spricht man bei
Einatmung von reizenden, giftigen und/oder heißen Gasen,
was eine Schädigung der Atemwege und der Lunge zur Folge
hat.}
%
Typische Leitsymptome sind:
\begin{itemize}
  \item Z.\,n. Gasexposition
  \item Angebrannte Gesichtshaare
  \item Geschwollene, rußige Lippen
  \item Rußpartikel im Auswurf
  \item Husten
  \item Heiserkeit
  \item Stridor
  \item Atemnot
\end{itemize}
Das volle Ausmaß der Schädigung kann unmittelbar oder
innerhalb von Stunden erreicht werden. Eine engmaschige
Beobachtung des Patienten ist daher sehr wichtig.
}
{%
\begin{itemize}
  \item Z.\,n. Gasexposition
  \item Angebrannte Gesichtshaare
  \item Geschwollene, rußige Lippen
  \item Rußpartikel im Auswurf
  \item Husten
  \item Heiserkeit
  \item Stridor
  \item Atemnot
\end{itemize}

Volles Ausmaß oft erst nach Stunden!
}





% M %%%%%%%%%%%%%%%%%%%%%%%%%%%%%%%%%%%%%%%%%%%%%%%%%%% M %
% Inhalationstrauma
\ProcedurePrintDefault{TT49091C}









%%%%%%%%%%%%%%%%%%%%%%%%%%%%%%%%%%%%%%%%%%%%%%%%%%%%%%%%%%%
%%%%%%%%%%%%%%%%%%%%%%%%%%%%%%%%%%%%%%%%%%%%%%%%%%%%%%%%%%%
\clearpage
\subsection{Verletzungen durch chemische Stoffe -- \enquote{Chemische Verbrennungen}}
\label{topic:veraetzung-trauma}%



Vergiftungen durch Einnahme von Säuren und Laugen werden 
unter \FullPrettyRef{topic:veraetzung-intoxikation}, besprochen. 

\Layout{60}{\TxBeschreibung}
{%
Schädigungen der Haut durch Säuren, Laugen oder andere
schädliche Stoffe weisen -- trotz der unterschiedlichen
Entstehung -- Gemeinsamkeiten mit Verbrennungsschäden auf.
Daher werden diese Verletzungen oft auch als
\emph{\enquote{chemische Verbrennung}} bezeichnet, auch wenn meist
keine Hitze im Spiel ist.

Genau wie bei der Verbrennung kommt es zu einem Ausfall der
Schutzfunktion der Haut, oft zu Schmerzen und zu
gefährlichen Schockzuständen. Auch die Einschätzung der
Ausdehnung der Verwundung erfolgt wie bei der Verbrennung
mit der Neuner-Regel o.\,ä.

Eine Besonderheit stellt jedoch der verantwortliche Stoff
dar, mit welchem u.\,U. der Patient, dessen Kleidung und
evtl. auch die Umgebung kontaminiert ist. Von diesem Stoff
kann sowohl für den Patienten, als auch für den Helfer eine
Gefahr ausgehen. Auf den Selbstschutz und die
\EE{Dekontamination} muss großer Wert gelegt werden.
}
{
\begin{itemize}
  \item Ausfall der Schutzfunktion der Haut
  \item Schmerzen
  \item gefährliche Schockzustände
  \item Gefahr für den Helfer
  \item Selbstschutz, Dekontamination!!!
\end{itemize}
}
{./images/trauma/notbrause2.jpg}
{}
{}

\begin{Literary}
\fullcite{RD200811Hoffmann}
\end{Literary}

% M %%%%%%%%%%%%%%%%%%%%%%%%%%%%%%%%%%%%%%%%%%%%%%%%%%% M %
% Verletzungen mit chemischen Substanzen
\ProcedurePrintDefault{TT30041C}









%%%%%%%%%%%%%%%%%%%%%%%%%%%%%%%%%%%%%%%%%%%%%%%%%%%%%%%%%%%
%%%%%%%%%%%%%%%%%%%%%%%%%%%%%%%%%%%%%%%%%%%%%%%%%%%%%%%%%%%
%%%%%%%%%%%%%%%%%%%%%%%%%%%%%%%%%%%%%%%%%%%%%%%%%%%%%%%%%%%
\clearpage % TEMP temp clearpage
\section{Erfrierungen}

\label{topic:erfrierung}%


\Text{\TxBeschreibung}
{%
\DefinitionInPara{Eine \T{Erfrierung} 
(\Z{\Lang{Lat.} \I{Congelatio} (\I{Congelat.})}) 
ist eine Gewebsnekrose, die durch lokale 
Durchblutungsstörungen aufgrund von  Kälteeinwirkung (Eis,
Wind, Nässe,  Kälte)
entstanden ist.}
%
}
Besonders gefährdet sind körperferne Körperteile wie Zehen,
Finger, Ohren oder Nase.
Die Einteilung erfolgt in drei \EE{Stadien} (siehe
\prettyref{tab:Erfrierungen}), diese haben am Notfallort
keine entscheidende Bedeutung für die zu setzenden
Maßnahmen. Den Schweregrad einer Erfrierung kann man erst
nach Tagen voll einschätzen. Deshalb ist es wichtig, die
Warnsymptome ernst zu nehmen!


\begin{table}[H]
\begin{center}
\Captionof{table}{Stadien der Erfrierung
\label{tab:Erfrierungen}} 
\begin{tabulary}{\linewidth}[H]{CLC}\addlinespace\toprule\addlinespace
\noindent\TabHeading{Stadien}
 &
\noindent\TabHeading{Symptome}
 &
\noindent~
\\\addlinespace\midrule\addlinespace
1
 &
Blässe, Empfindungsstörungen
 &
\\\addlinespace
2
 &
Blasenbildung, Schmerzen
 &
\\\addlinespace
3
 &
grau/schwarze Marmorierung (Nekrose), Gefühllosigkeit
\\\addlinespace\bottomrule
\end{tabulary}
\end{center}
\end{table}


% M %%%%%%%%%%%%%%%%%%%%%%%%%%%%%%%%%%%%%%%%%%%%%%%%%%% M %
% Erfrierungen
\ProcedurePrintDefault{TT35071C}



%%%%%%%%%%%%%%%%%%%%%%%%%%%%%%%%%%%%%%%%%%%%%%%%%%%%%%%%%%%
%%%%%%%%%%%%%%%%%%%%%%%%%%%%%%%%%%%%%%%%%%%%%%%%%%%%%%%%%%%
%%%%%%%%%%%%%%%%%%%%%%%%%%%%%%%%%%%%%%%%%%%%%%%%%%%%%%%%%%%
\clearpage % TEMP temp clearpage
\section{Unfälle durch Strom- und Blitzschlag}

\marginpar{\includegraphics[width=\linewidth]{./images/electricity_japan-00800.jpg}
\small\emph{Foto: Fina}}

Stromunfälle sind zwar selten, stellen jedoch für die Helfer eine große
Herausforderung dar. Der Grund hierfür ist, dass bereits im Szeneüberblick die
Gefahr eines möglichen weiteren Stromschlages erkannt werden muss! Erster
Anhaltspunkt ist daher die Situation am Einsatzort. Typisch sind:
\begin{itemize}
  \item Spielende Jugendliche oder Arbeiter auf Eisenbahnwaggons
  \item LKW-Hebekräne oder Bagger, die in Freileitungen geraten sind
  \item Unachtsamer Umgang mit Elektrogeräten im Badezimmer
  \item Unsachgemäße Reparatur von Elektrogeräten
  \item Arbeitsunfälle (Elektriker, Heimwerker)
  \item Verkehrsunfall mit herabhängenden Stromleitungen
  \item Blitzunfall
\end{itemize}
Die erste Maßnahme muss daher immer dem Selbstschutz gelten:
\begin{itemize}
  \item Strom abschalten (lassen)! ${\rightarrow}$ Sicherungen, Abziehen des
  Netzkabels, NOT-AUS-Taste, Abschalten ganzer Anlagen durch den Fachmann, \ldots
  \item Sicherheitsabstand einhalten! Bei Hochspannungsleitungen sind dies
  mindestens 5\,m. \cite{BoehmerTaschRett6}
  \item Bahnanlagen sperren und sichern lassen!
  \item Strom gegen Wiedereinschalten sichern (lassen)!
\end{itemize}
Die Einschätzung über das Ausmaß der Verletzung bzw. der vitalen Bedrohung ist
bei Elektrounfällen sehr schwierig. Dies hat mehrere Gründe:
\begin{enumerate}
  \item Die Wirkung des Stromes im menschlichen Körper hängt von mehreren
  Faktoren ab: Höhe der elektrischen Spannung\footnote{Die elektrische Spannung
  wird in Volt (V) gemessen. Typische Werte sind 12\,V bei einer Autobatterie,
  230\,V an einer Haushaltssteckdose, 15\,000\,V bei Eisenbahn-Oberleitungen
  und bis zu 380\,000\,V bei Hochspannungsleitungen. Die elektrische Spannung
  ist die treibende Kraft hinter dem Stromfluss. Die tatsächliche Stärke des
  Stromes hängt jedoch auch vom elektrischen Widerstand des durchströmten
  Materials ab.}, Stärke des elektrischen Stromes\footnote{Die Stärke des
  elektrischen Stromes wird in Ampere (A) gemessen.}, der
  Stromart (Gleich- oder Wechselstrom), dem Weg des Stromes
  durch den Körper, dem Kontaktwiderstand (feuchte Hände,
  Wasserlacke, \ldots) und der Einwirkzeit. Meistens sind
  die Zahlenwerte dieser Faktoren im Einsatzfall jedoch
  nicht bekannt, sodass eine theoretische Abschätzung der
  Gefährung praktisch nicht durchführbar ist. 
  \item Die \I{Strommarken} (Verbrennungen an den Ein- und
  Austrittsstellen) sind i.\,d.\,R. klein, sodass die Schädung im Körperinneren nicht abgeschätzt
  werden kann. Im Zweifelsfall muss man mit ausgedehnten inneren Verletzungen
  rechnen.
\end{enumerate}
Die Orientierung an Symptomen ist für die Einschätzung der vitalen Bedrohung
daher oft sinnvoller.

\TextSlide{\TxSymptome}
{%
\begin{itemize}
  \item Szeneüberblick: Unfallmechanismus und Unfallhergang
  \item Reizeffekte (Muskel-/ Nervenerregung):
  \begin{itemize}
    \item Verkrampfung ({\quotedblbase}kann den stromführenden Teil nicht
    loslassen{\textquotedblleft})
    \item Sensibilitätsstörungen
    \item ZNS ${\rightarrow}$ Bewußtseinsstörung
    \item Herzrhythmusstörungen (Kammerflimmern, evtl.  als Spätfolge!)
  \end{itemize}
  \item Bewusstseinsstörungen bis Krampfanfälle oder Bewusstlosigkeit
  \item Lähmungen (vorübergehend)
  \item Thermische Schäden durch Wärmeentstehung
  \begin{itemize}
    \item Strommarken (Ein-/ Austrittsverbrennung)
    \item Lichtbogen: oberflächliche und tiefe  Verbrennungen
    \item Blitzfiguren an der Haut
  \end{itemize}
\end{itemize}

% Stromunfälle kommen typischerweise in folgenden Situationen vor:
% \begin{itemize}
%   \item Spielende Jugendliche oder Arbeiter auf Eisenbahnwaggons
%   \item LKW-Hebekräne oder Bagger geraten in Freileitungen
%   \item Unachtsamer Umgang mit Elektrogeräten im Badezimmer
%   \item Unsachgemäße Reparatur von Elektrogeräten
%   \item Arbeitsunfälle (Elektriker, Heimwerker)
%   \item Tragen von  Metallteilen in der Nähe von Hochspannungsleitungen
% \end{itemize}


% \TextSlide{Allgemeines zum elektrischen Strom und seiner
% Wirkung auf den Körper}
% {%
% Die Gefährdung bei Stromunfällen hängt von verschiedenen Faktoren ab:
% \begin{description}
%   \item[Elektrische Spannung:] Die elektrische Spannung ist
%   die Ursache für den elektrischen Stromfluss und wird in
%   der Einheit \T{Volt} (V) gemessen.\Z{Man
%   unterscheidet zwischen Hochspannung (über 1000\,V) und
%   Niederspannung (unter 1000\,V).} Grundsätzlich gilt: Je
%   höher die Spannung ist, desto gefährlicher ist es.\ZFootnote{An
%   einer Autobatterie liegt eine Spannung von 12 V. An einer Steckdose liegt eine Spannung von 230 V, in
%     Industriebetrieben wird zusätzlich oft eine Spannung von
%     400 V verwendet. Diese Spannungen sind auf jeden Fall
%     lebensgefährlich. Ein Defibrillator schockt mit etwa
%     1000\,V (!), ein externer Herzschrittmacher arbeitet mit
%     bis zu 400\,V. Über
%     Hochspannungsleitungen wird die elektrische Energie mit
%     einer Spannung von 110\,000\,V (= 110\,kV) bis
%     380\,000\,V (= 380\,kV) übertragen. An den Oberleitungen
%     der Eisenbahn liegen 15\,000\,V (= 15\,kV). Bei
%     Gewitterblitzen treten Spannungen von 10-30 Millionen
%     Volt (10-30 MV) auf.} Wie gefährlich eine
%   Spannung aber tatsächlich ist hängt nicht nur von der Höhe
%   der Spannung selbst, sondern wesentlich auch von der Stärke des
%   elektrischen Stroms, der von ihr hervorgerufen wird,
%   ab! So kann z.\,B. die 12\,V-Spannung einer
%   Modelleisenbahn völlig ungefährlich sein, gibt es jedoch im Auto bei einer
%   (großen) 12\,V-Batterie einen Kurzschluss, kann ein
%   lebensgefährlich großer Strom fließen. Meistens kann am
%   Notfallort jedoch nur die Höhe der Spannung abgeschätzt
%   werden.
% %   \begin{itemize}
% %     \item \T{Niederspannung}: Von Niederspannung spricht man
% %     bei Spannungen bis 1000\,V. \item \T{Hochspannung} ist eine
% %     Spannung über 1000\,V. 
% %   \end{itemize}
%   \item[Stromstärke:] Der elektrische 
%   Strom  wird durch die elektrische
%   \T{Stromstärke} beschrieben, welche mit der  Einheit
%   \T{Ampere} (A) gemessen wird.\ZFootnote{Elektrischer Strom ist der Fluss von
%   Elektronen.  Die treibende Kraft hinter dem Stromfluss ist
%   die elektrische Spannung.} Kommt es  zu einem
%   Stromunfall ist die Schädigung im Körper von der Höhe der
%   Stromstärke  abhängig.\ZFootnote{Bei geringen Stromstärken
%   \E{unter 1 Milliampere (mA)}  ist  nur ein elektrischer
%   Schlag spürbar. Bei Stromstärken \E{ab ca.  10 Milliampere
%   (mA)} kommt es zu Schmerzempfindungen und 
%   Muskelkontraktionen \cite{Jaros1}, ab \EE{15-25 mA} kommt
%   es aufgrund der Muskelkontraktionen zu  krampfhaftem
%   Festhalten der berührten Stromquelle
%   \cite{BoehmerTaschRett6}.  Stromstärken ab
%   \E{25-50\,mA} sind \E{lebensgefährlich}!~\cite{BoehmerTaschRett6}
%   Ab \E{80\,mA} kommt es zu Atemlähmung und Kammerflimmern
%   \cite{Jaros1}, Stärken über \E{3\,A} führen zusätzlich zu
%   Verbrennungen, Verkochungen und Verkohlungen \cite{LPN3}. 
%   Die Werte beziehen sich auf Wechselstrom, die Grenzwerte
%   für Gleichstrom sind höher. Da oft die Stromart nicht
%   bekannt ist, ist es sinnvoll im Notfall von der
%   gefährlicheren  Stromart auszugehen!} Stromstärken ab
%   25\,mA gelten als potentiell tötlich! Im Allgemeinen wird
%   am Notfallort die Stromstärke unbekannt sein.
%   \item[Widerstand:]
%   Jeder Körper und
%   jeder Stoff setzt dem elektrischen Stromfluss einen
%   Widerstand entgegen. Dieser elektrische Widerstand wird in
%   \T{Ohm} ($\Omega$) gemessen. Je größer der Widerstand ist,
%   desto kleiner ist der elektrische Strom. \EE{Feuchte Haut hat einen deutlich geringeren Widerstand als
%   trockene Haut.} Der Widerstand von geschädigter
%   (z.\,B. verbrannter) Haut ist noch
%   geringer.\ZFootnote{Gerät man mit feuchter Haut in die
%   Steckdose ist der Widerstand so gering, dass die
%   Stromstärke lebensbedrohlich hoch wird.}
%   \item[Einwirkdauer:] Je länger der Stromfluss gedauert
%   hat, desto größer ist die Schädigung.
%   \item[Stromweg:] Die Schädigung im Körper hängt auch vom
%   Weg des Stromes durch den Körper ab. Fließt der Strom
%   z.\,B. von einer Hand zur anderen Hand, so fließt er
%   längs über das Herz und ist somit gefährlicher,
%   als wenn der Strom von einer Hand zum
%   Fuß fließt. Um den Stromweg heraus zu finden, müssen
%   Eintritts- und Austrittsstellen (Strommarken) am Körper
%   gesucht werden.
%   \item[Stromart:] Es gibt Wechselstrom und
%   Gleichstrom. An den Steckdosen liegt
%   Wechselspannung an, aus Batterien fließt Gleichstrom.\ZFootnote{Bei monophasischen
%   Defibrillatoren wird mit Gleichstrom gearbeitet, biphasische Defibrillatoren polen die Stromrichtung einmal um (Wechselstrom).}
%   Wechselstrom ist gefährlicher als Gleichstrom. Bei Kontakt
%   mit Wechselstrom ist die Wahrscheinlichkeit für
%   Kammerflimmern größer als bei Gleichstrom
%   \cite{RedelsteinerHRNS1}.
% \end{description}
% 
% 
% %\T{Niederspannung}: bis 1000\,Volt (Steckdose:  230 Volt)
% %\T{Hochspannung}: über 1000\,Volt, Stromstärke
% %{\textgreater}  3 Ampere (Überlandleitungen) 
% 
% }{
% \begin{itemize}
%   \item Je höher die Spannung, desto gefährlicher!
%   \item Je höher die Stromstärke, desto gefährlicher! Stromstärken ab 15 mA potentiell tödlich!
%   \item Je geringer der Widerstand, desto gefährlicher!
%   \item Je länger die Einwirkdauer, desto gefährlicher!
%   \item Stromweg längs durch das Herz gefährlicher als
%   anderswo!
%   \item Wechselstrom ist gefährlicher als Gleichstrom!
% \end{itemize}
% }{}{}{}

\AuthNotes{\#9: EMHOFER: Bilder von verschiedenen
Steckdosen, Kabel und Hinweisschildern einfügen}





% \TextSlide{\TxSymptome}
% {%
% Die Verletzungsmuster und Symptome variieren je nach
% Spannung (Volt) und Stromstärke (Ampere). Sie können daher
% sehr unterschiedlich sein:
% 
% \begin{itemize}
%   \item Unfallmechanismus und Unfallhergang
%   \item Reizeffekte (Muskel-/ Nervenerregung):
%   \begin{itemize}
%     \item Verkrampfung ({\quotedblbase}kann den stromführenden Teil nicht
%     loslassen{\textquotedblleft})
%     \item Sensibilitätsstörungen
%     \item ZNS ${\rightarrow}$ Bewußtseinsstörung
%     \item Herzrhythmusstörungen (Kammerflimmern, evtl.  als Spätfolge!)
%   \end{itemize}
%   \item Bewusstseinsstörungen bis Krampfanfälle oder Bewusstlosigkeit
%   \item Lähmungen (vorübergehend)
%   \item Thermische Schäden durch Wärmeentstehung
%   \begin{itemize}
%     \item Strommarken (Ein-/ Austrittsverbrennung)
%     \item Lichtbogen: oberflächliche und tiefe  Verbrennungen
%     \item Blitzfiguren an der Haut
%   \end{itemize}
% \end{itemize}
%\begin{itemize}
%    \item Gleichstrom
%    \item Wechselstrom ${\rightarrow}$ eher  Herzrhythmusstörungen
%    \item Patient  kann Spannungsquelle evtl. nicht  loslassen (Krampf der
%    Handmuskulatur)
%    \item Auch Defi stellt Strom-Gefahrenquelle  dar!!!
%    \item Blitzunfall
%    \begin{itemize}
%        \item Mehrere Millionen (!) Volt
%        \item Bis 100.000 Ampere
%    \end{itemize}
%\end{itemize}

%\paragraph{\TxSymptome}
%\begin{itemize}
%    \item Bewußtlosigkeit
%    \item Lähmungen (vorübergehend)
%    \item Blitzfiguren auf der Haut
%    \item Herzrhythmusstörungen möglich
%\end{itemize}

\Cave{Das eigentliche Ausmaß der Verletzung ist oft
\EE{{\quotedblbase}unsichtbar{\textquotedblleft}}, da der
Strom im Körperinneren -- also {\quotedblbase}unterirdisch{\textquotedblleft}
-- viel Schaden anrichtet. }
}
{
\begin{itemize}
  \item Reizeffekte (Muskel-/ Nervenerregung)
  \item Bewusstseinsstörungen bis Krampfanfälle oder Bewusstlosigkeit
  \item Thermische Schäden, z.\,B. Strommarken
\end{itemize}
}{}{}{}


\TextSlide{Spezielle Mechanismen und Situationen}{%

\begin{description}
  \item[Stromüberlandleitungen (Hochspannungsleitungen)]
  Bei Überlandleitungen ist mit Hochspannung (380\,000 Volt
  (380\,kV)) zu rechnen. Ist eine Hochspannungsleitung
  gerissen und berührt den Boden, so ist besondere Vorsicht
  geboten! Da auch das Erdreich einen elektrischen Widerstand hat, gibt es in der
  Nähe der gerissenen Leitung im Boden einen
  Spannungsabfall. Bewegt man sich in großen Schritten auf
  das Kabel zu oder vom Kabel weg, so tritt zwischen den
  Füßen eine derart hohe Spannung auf, dass es zu einem
  Stromfluss über den menschlichen Körper von einem Fuß zum
  anderen Fuß kommen kann (\EE{Schrittspannung}). Bei
  gerissenen Hochspannungsleitungen darf man daher keine großen
  Schritte machen, oder stromleitende Gegenstände (z.\,B.
  Matalle) berühren! \prettyref{topic:gefahrenzonen}
  \item[Bahnanlagen]
  Die Eisenbahn arbeitet mit einer Oberleitungsspannung von
  15\,000\,V, auch bei gesenktem Stromabnehmer oder
  Dieselfahrzeugen kann durch Kondensatorladung eine Spannung
  von bis zu 3\,000\,V anliegen. \prettyref{topic:gefahrenzonen}
  
  Im U-Bahnbereich werden
  Spannungen um die 700\,V eingesetzt. \cite{U61989}
\end{description}

}{%
\begin{itemize}
  \item Stromüberlandleitungen ${\rightarrow}$ Schrittspannung
  \item Bahnanlagen
\end{itemize}
}{}{}{}

\TextSlide{Blitzschlag}{%
Bei Gewitterblitzen treten Spannungen von Millionen von Volt
und Stromstärken von bis zu 300\,000 Ampere auf. Trotzdem
gibt es auch bei Blitzunfällen immer wieder Überlebende, da
die Einwirkdauer extrem kurz (wenige Millisekunden) ist.
Typische Symptome sind
\begin{itemize}
  \item Bewusstlosigkeit
  \item Vorübergehende Lähmung
  \item Blitzfiguren auf der Haut
  \item Herzrhythmusstörungen
\end{itemize}
}{%
\begin{itemize}
  \item extrem kurze Einwirkdauer
  \item typische Symptome
\end{itemize}
}{}{}{}


% M %%%%%%%%%%%%%%%%%%%%%%%%%%%%%%%%%%%%%%%%%%%%%%%%%%% M %
% Stromunfälle
\ProcedurePrintDefault{TT75040C}

\section*{Weiterführende Literatur}
\noindent\SymbolLiteratur{} 
\cite{PHTLS6,
emergency9,
OHCSP7,
Helfen2008,
ITLSde6,
DRAn3,
Patzelt2005,
Ficklscherer2005,
Sachsenweger2003,
SiewertChirurgie8,
ThierbachInterhospitaltransfer1,
RueckerBildatlas1,
ArbeitstechnikenAZ1,
ScholzNotfallmedizin2}
% \nocite{PHTLS6}
% \nocite{emergency9}
% \nocite{OHCSP7}
% \nocite{Helfen2008}
% \nocite{DRAn3}
% \nocite{Patzelt2005}
% \nocite{Ficklscherer2005}
% \nocite{Sachsenweger2003}
% \nocite{SiewertChirurgie8}
% \nocite{ThierbachInterhospitaltransfer1}
% \nocite{RueckerBildatlas1}
% \nocite{ArbeitstechnikenAZ1}
% \nocite{ScholzNotfallmedizin2}



\clearpage
\ChapterBibliography
\ChapterEnd
